% Options for packages loaded elsewhere
% Options for packages loaded elsewhere
\PassOptionsToPackage{unicode}{hyperref}
\PassOptionsToPackage{hyphens}{url}
\PassOptionsToPackage{dvipsnames,svgnames,x11names}{xcolor}
%
\documentclass[
  spanish,
  a4paper,
  oneside]{scrbook}
\usepackage{xcolor}
\usepackage[top=25mm,bottom=25mm,left=30mm,right=20mm]{geometry}
\usepackage{amsmath,amssymb}
\setcounter{secnumdepth}{5}
\usepackage{iftex}
\ifPDFTeX
  \usepackage[T1]{fontenc}
  \usepackage[utf8]{inputenc}
  \usepackage{textcomp} % provide euro and other symbols
\else % if luatex or xetex
  \usepackage{unicode-math} % this also loads fontspec
  \defaultfontfeatures{Scale=MatchLowercase}
  \defaultfontfeatures[\rmfamily]{Ligatures=TeX,Scale=1}
\fi
\usepackage{lmodern}
\ifPDFTeX\else
  % xetex/luatex font selection
\fi
% Use upquote if available, for straight quotes in verbatim environments
\IfFileExists{upquote.sty}{\usepackage{upquote}}{}
\IfFileExists{microtype.sty}{% use microtype if available
  \usepackage[]{microtype}
  \UseMicrotypeSet[protrusion]{basicmath} % disable protrusion for tt fonts
}{}
\makeatletter
\@ifundefined{KOMAClassName}{% if non-KOMA class
  \IfFileExists{parskip.sty}{%
    \usepackage{parskip}
  }{% else
    \setlength{\parindent}{0pt}
    \setlength{\parskip}{6pt plus 2pt minus 1pt}}
}{% if KOMA class
  \KOMAoptions{parskip=half}}
\makeatother
% Make \paragraph and \subparagraph free-standing
\makeatletter
\ifx\paragraph\undefined\else
  \let\oldparagraph\paragraph
  \renewcommand{\paragraph}{
    \@ifstar
      \xxxParagraphStar
      \xxxParagraphNoStar
  }
  \newcommand{\xxxParagraphStar}[1]{\oldparagraph*{#1}\mbox{}}
  \newcommand{\xxxParagraphNoStar}[1]{\oldparagraph{#1}\mbox{}}
\fi
\ifx\subparagraph\undefined\else
  \let\oldsubparagraph\subparagraph
  \renewcommand{\subparagraph}{
    \@ifstar
      \xxxSubParagraphStar
      \xxxSubParagraphNoStar
  }
  \newcommand{\xxxSubParagraphStar}[1]{\oldsubparagraph*{#1}\mbox{}}
  \newcommand{\xxxSubParagraphNoStar}[1]{\oldsubparagraph{#1}\mbox{}}
\fi
\makeatother

\usepackage{color}
\usepackage{fancyvrb}
\newcommand{\VerbBar}{|}
\newcommand{\VERB}{\Verb[commandchars=\\\{\}]}
\DefineVerbatimEnvironment{Highlighting}{Verbatim}{commandchars=\\\{\}}
% Add ',fontsize=\small' for more characters per line
\usepackage{framed}
\definecolor{shadecolor}{RGB}{241,243,245}
\newenvironment{Shaded}{\begin{snugshade}}{\end{snugshade}}
\newcommand{\AlertTok}[1]{\textcolor[rgb]{0.68,0.00,0.00}{#1}}
\newcommand{\AnnotationTok}[1]{\textcolor[rgb]{0.37,0.37,0.37}{#1}}
\newcommand{\AttributeTok}[1]{\textcolor[rgb]{0.40,0.45,0.13}{#1}}
\newcommand{\BaseNTok}[1]{\textcolor[rgb]{0.68,0.00,0.00}{#1}}
\newcommand{\BuiltInTok}[1]{\textcolor[rgb]{0.00,0.23,0.31}{#1}}
\newcommand{\CharTok}[1]{\textcolor[rgb]{0.13,0.47,0.30}{#1}}
\newcommand{\CommentTok}[1]{\textcolor[rgb]{0.37,0.37,0.37}{#1}}
\newcommand{\CommentVarTok}[1]{\textcolor[rgb]{0.37,0.37,0.37}{\textit{#1}}}
\newcommand{\ConstantTok}[1]{\textcolor[rgb]{0.56,0.35,0.01}{#1}}
\newcommand{\ControlFlowTok}[1]{\textcolor[rgb]{0.00,0.23,0.31}{\textbf{#1}}}
\newcommand{\DataTypeTok}[1]{\textcolor[rgb]{0.68,0.00,0.00}{#1}}
\newcommand{\DecValTok}[1]{\textcolor[rgb]{0.68,0.00,0.00}{#1}}
\newcommand{\DocumentationTok}[1]{\textcolor[rgb]{0.37,0.37,0.37}{\textit{#1}}}
\newcommand{\ErrorTok}[1]{\textcolor[rgb]{0.68,0.00,0.00}{#1}}
\newcommand{\ExtensionTok}[1]{\textcolor[rgb]{0.00,0.23,0.31}{#1}}
\newcommand{\FloatTok}[1]{\textcolor[rgb]{0.68,0.00,0.00}{#1}}
\newcommand{\FunctionTok}[1]{\textcolor[rgb]{0.28,0.35,0.67}{#1}}
\newcommand{\ImportTok}[1]{\textcolor[rgb]{0.00,0.46,0.62}{#1}}
\newcommand{\InformationTok}[1]{\textcolor[rgb]{0.37,0.37,0.37}{#1}}
\newcommand{\KeywordTok}[1]{\textcolor[rgb]{0.00,0.23,0.31}{\textbf{#1}}}
\newcommand{\NormalTok}[1]{\textcolor[rgb]{0.00,0.23,0.31}{#1}}
\newcommand{\OperatorTok}[1]{\textcolor[rgb]{0.37,0.37,0.37}{#1}}
\newcommand{\OtherTok}[1]{\textcolor[rgb]{0.00,0.23,0.31}{#1}}
\newcommand{\PreprocessorTok}[1]{\textcolor[rgb]{0.68,0.00,0.00}{#1}}
\newcommand{\RegionMarkerTok}[1]{\textcolor[rgb]{0.00,0.23,0.31}{#1}}
\newcommand{\SpecialCharTok}[1]{\textcolor[rgb]{0.37,0.37,0.37}{#1}}
\newcommand{\SpecialStringTok}[1]{\textcolor[rgb]{0.13,0.47,0.30}{#1}}
\newcommand{\StringTok}[1]{\textcolor[rgb]{0.13,0.47,0.30}{#1}}
\newcommand{\VariableTok}[1]{\textcolor[rgb]{0.07,0.07,0.07}{#1}}
\newcommand{\VerbatimStringTok}[1]{\textcolor[rgb]{0.13,0.47,0.30}{#1}}
\newcommand{\WarningTok}[1]{\textcolor[rgb]{0.37,0.37,0.37}{\textit{#1}}}

\usepackage{longtable,booktabs,array}
\newcounter{none} % for unnumbered tables
\usepackage{calc} % for calculating minipage widths
% Correct order of tables after \paragraph or \subparagraph
\usepackage{etoolbox}
\makeatletter
\patchcmd\longtable{\par}{\if@noskipsec\mbox{}\fi\par}{}{}
\makeatother
% Allow footnotes in longtable head/foot
\IfFileExists{footnotehyper.sty}{\usepackage{footnotehyper}}{\usepackage{footnote}}
\makesavenoteenv{longtable}
\usepackage{graphicx}
\makeatletter
\newsavebox\pandoc@box
\newcommand*\pandocbounded[1]{% scales image to fit in text height/width
  \sbox\pandoc@box{#1}%
  \Gscale@div\@tempa{\textheight}{\dimexpr\ht\pandoc@box+\dp\pandoc@box\relax}%
  \Gscale@div\@tempb{\linewidth}{\wd\pandoc@box}%
  \ifdim\@tempb\p@<\@tempa\p@\let\@tempa\@tempb\fi% select the smaller of both
  \ifdim\@tempa\p@<\p@\scalebox{\@tempa}{\usebox\pandoc@box}%
  \else\usebox{\pandoc@box}%
  \fi%
}
% Set default figure placement to htbp
\def\fps@figure{htbp}
\makeatother


% definitions for citeproc citations
\NewDocumentCommand\citeproctext{}{}
\NewDocumentCommand\citeproc{mm}{%
  \begingroup\def\citeproctext{#2}\cite{#1}\endgroup}
\makeatletter
 % allow citations to break across lines
 \let\@cite@ofmt\@firstofone
 % avoid brackets around text for \cite:
 \def\@biblabel#1{}
 \def\@cite#1#2{{#1\if@tempswa , #2\fi}}
\makeatother
\newlength{\cslhangindent}
\setlength{\cslhangindent}{1.5em}
\newlength{\csllabelwidth}
\setlength{\csllabelwidth}{3em}
\newenvironment{CSLReferences}[2] % #1 hanging-indent, #2 entry-spacing
 {\begin{list}{}{%
  \setlength{\itemindent}{0pt}
  \setlength{\leftmargin}{0pt}
  \setlength{\parsep}{0pt}
  % turn on hanging indent if param 1 is 1
  \ifodd #1
   \setlength{\leftmargin}{\cslhangindent}
   \setlength{\itemindent}{-1\cslhangindent}
  \fi
  % set entry spacing
  \setlength{\itemsep}{#2\baselineskip}}}
 {\end{list}}
\usepackage{calc}
\newcommand{\CSLBlock}[1]{\hfill\break\parbox[t]{\linewidth}{\strut\ignorespaces#1\strut}}
\newcommand{\CSLLeftMargin}[1]{\parbox[t]{\csllabelwidth}{\strut#1\strut}}
\newcommand{\CSLRightInline}[1]{\parbox[t]{\linewidth - \csllabelwidth}{\strut#1\strut}}
\newcommand{\CSLIndent}[1]{\hspace{\cslhangindent}#1}

\ifLuaTeX
\usepackage[bidi=basic,shorthands=off]{babel}
\else
\usepackage[bidi=default,shorthands=off]{babel}
\fi
\ifLuaTeX
  \usepackage{selnolig} % disable illegal ligatures
\fi


\setlength{\emergencystretch}{3em} % prevent overfull lines

\providecommand{\tightlist}{%
  \setlength{\itemsep}{0pt}\setlength{\parskip}{0pt}}



 


\AtBeginDocument{\hypersetup{linkcolor=black}}
% ---------- Paquetes tipograficos y microajustes ----------
\usepackage{microtype}        % mejor interletrado y justificado
\usepackage{csquotes}         % comillas tipograficas
\usepackage{iftex}
\ifPDFTeX\else
  % Seleccion de fuentes solo cuando XeTeX/LuaTeX esta disponible.
  \IfFontExistsTF{TeX Gyre Termes}{
    \setmainfont{TeX Gyre Termes}
  }{
    \IfFontExistsTF{Latin Modern Roman}{\setmainfont{Latin Modern Roman}}{}
  }
  \IfFontExistsTF{TeX Gyre Heros}{
    \setsansfont{TeX Gyre Heros}
  }{
    \IfFontExistsTF{Latin Modern Sans}{\setsansfont{Latin Modern Sans}}{}
  }
  \IfFontExistsTF{TeX Gyre Cursor}{
    \setmonofont{TeX Gyre Cursor}
  }{
    \IfFontExistsTF{Latin Modern Mono}{\setmonofont{Latin Modern Mono}}{}
  }
\fi
\usepackage{enumitem}         % listas compactas
\setlist{itemsep=.2em, topsep=.2em}

% ---------- KOMA-Script: estilo de titulos y espaciado ----------
\KOMAoptions{
  headings=big,
  parskip=half,
  fontsize=12pt,
  appendixprefix=true
}
\setlength{\parindent}{1.5em}
\setlength{\parskip}{0.6em}
\usepackage{indentfirst}
\usepackage{setspace}         % Control del interlineado del documento
\setstretch{1.5}

% ---------- Encabezados y pies (scrlayer-scrpage) ----------
\usepackage[automark,headsepline]{scrlayer-scrpage}
\clearpairofpagestyles
\automark[chapter]{chapter}
\ihead{\pagemark}
\ohead{\itshape\headmark}
\setheadsepline{0.4pt}
\renewcommand*{\chaptermarkformat}{}
\renewcommand*{\chapterpagestyle}{scrheadings}
\pagestyle{scrheadings}
\makeatletter
\let\ps@plain\ps@scrheadings
\makeatother
\cfoot{}

% ---------- Hipervinculos mas sobrios ----------
\usepackage{hyperref}
\hypersetup{
  colorlinks=true,
  linkcolor=black,
  citecolor=[rgb]{0.05,0.2,0.5},
  urlcolor=[rgb]{0.05,0.2,0.5},
  pdfauthor={\@author},
  pdftitle={\@title}
}

% ---------- Leyendas de figuras/tablas ----------
\usepackage[labelfont=bf,textfont=it]{caption}
\captionsetup{
  skip=8pt
}

% ---------- Tabla de contenidos ----------
\KOMAoptions{toc=graduated}
\RedeclareSectionCommand[tocnumwidth=3em]{chapter}
\RedeclareSectionCommand[tocindent=3.25em,tocnumwidth=2.8em]{section}
\RedeclareSectionCommand[tocindent=6.5em,tocnumwidth=2.5em]{subsection}
\makeatletter
\newcommand*{\tocdotfill}{\leavevmode\leaders\hbox to .6em{\hss.\hss}\hfill}
\makeatother
\RedeclareSectionCommand[toclinefill=\tocdotfill]{chapter}
\RedeclareSectionCommand[toclinefill=\tocdotfill]{section}
\RedeclareSectionCommand[toclinefill=\tocdotfill]{subsection}
\renewcommand*{\contentsname}{Tabla de contenido}
\setkomafont{chapterentry}{\normalfont}
\setkomafont{chapterentrypagenumber}{\normalfont}

% ---------- Viudas/Huerfanas y cortes de pagina ----------
\clubpenalty=10000
\widowpenalty=10000
\displaywidowpenalty=10000

% ---------- Entorno abstract para scrbook ----------
\providecommand{\abstractname}{Resumen}
\makeatletter
\@ifundefined{abstract}{
  \newenvironment{abstract}{
    \cleardoublepage
    \thispagestyle{plain}
    \null\vfill
    \begin{center}
      {\bfseries\Large \abstractname\par}
    \end{center}\vspace{1em}
    \begingroup
  }{
    \par\endgroup
    \vfill\null
    \cleardoublepage
  }
}{}
\makeatother

% ---------- Soporte de subtitulo desde YAML ----------
\makeatletter
\providecommand{\subtitle}[1]{\gdef\@subtitle{#1}}
\providecommand{\@subtitle}{}
\makeatother

\makeatletter
\newcommand{\facso@iflist}[4]{%
  \begingroup
    \def\addvspace##1{}%
    \def\numberline##1{##1}%
    \def\facso@target{#2}%
    \newcount\facso@entrycount
    \facso@entrycount=0\relax
    \long\def\contentsline##1##2##3{%
      \def\facso@this{##1}%
      \ifx\facso@this\facso@target
        \advance\facso@entrycount\@ne
      \fi
    }%
    \InputIfFileExists{\jobname.#1}{}{}%
  \endgroup
  \ifnum\facso@entrycount>0\relax
    #3%
  \else
    #4%
  \fi}
\renewcommand{\listoffigures}{%
  \facso@iflist{lof}{figure}{%
    \cleardoublepage
    \phantomsection
    \addchap*{Listado de Figuras}%
    \thispagestyle{scrheadings}%
    \markboth{Listado de Figuras}{Listado de Figuras}%
    \addcontentsline{toc}{chapter}{Listado de Figuras}%
    \@starttoc{lof}%
    \cleardoublepage
    \automark[chapter]{chapter}%
  }{}}
\renewcommand{\listoftables}{%
  \facso@iflist{lot}{table}{%
    \cleardoublepage
    \phantomsection
    \addchap*{Listado de Tablas}%
    \thispagestyle{scrheadings}%
    \markboth{Listado de Tablas}{Listado de Tablas}%
    \addcontentsline{toc}{chapter}{Listado de Tablas}%
    \@starttoc{lot}%
    \cleardoublepage
    \automark[chapter]{chapter}%
  }{}}
\makeatother

% --- Desactivar portada y abstract automaticos de Pandoc (PDF) ---
\AtBeginDocument{\let\maketitle\relax}
\renewenvironment{abstract}{}{}
\providecommand{\appendixname}{}
\providecommand{\appendixtocname}{}
\providecommand{\appendixpagename}{}
\renewcommand*{\appendixname}{Anexo}
\renewcommand*{\appendixtocname}{Anexos}
\renewcommand*{\appendixpagename}{Anexos}
\makeatletter
\@ifpackageloaded{bookmark}{}{\usepackage{bookmark}}
\makeatother
\makeatletter
\@ifpackageloaded{caption}{}{\usepackage{caption}}
\AtBeginDocument{%
\ifdefined\contentsname
  \renewcommand*\contentsname{Tabla de contenidos}
\else
  \newcommand\contentsname{Tabla de contenidos}
\fi
\ifdefined\listfigurename
  \renewcommand*\listfigurename{Listado de Figuras}
\else
  \newcommand\listfigurename{Listado de Figuras}
\fi
\ifdefined\listtablename
  \renewcommand*\listtablename{Listado de Tablas}
\else
  \newcommand\listtablename{Listado de Tablas}
\fi
\ifdefined\figurename
  \renewcommand*\figurename{Lista de figuras}
\else
  \newcommand\figurename{Lista de figuras}
\fi
\ifdefined\tablename
  \renewcommand*\tablename{Lista de tablas}
\else
  \newcommand\tablename{Lista de tablas}
\fi
}
\@ifpackageloaded{float}{}{\usepackage{float}}
\floatstyle{ruled}
\@ifundefined{c@chapter}{\newfloat{codelisting}{h}{lop}}{\newfloat{codelisting}{h}{lop}[chapter]}
\floatname{codelisting}{Listado}
\newcommand*\listoflistings{\listof{codelisting}{Listado de Listados}}
\makeatother
\makeatletter
\makeatother
\makeatletter
\@ifpackageloaded{caption}{}{\usepackage{caption}}
\@ifpackageloaded{subcaption}{}{\usepackage{subcaption}}
\makeatother
\usepackage{bookmark}
\IfFileExists{xurl.sty}{\usepackage{xurl}}{} % add URL line breaks if available
\urlstyle{same}
\hypersetup{
  pdftitle={Percepciones meritocráticas: percepciones de esfuerzo y talento en la distribución de recompensas sus factores asociados},
  pdfauthor={Bastián Javier Martin Gerdes},
  pdflang={es},
  colorlinks=true,
  linkcolor={Maroon},
  filecolor={Maroon},
  citecolor={Blue},
  urlcolor={Blue},
  pdfcreator={LaTeX via pandoc}}


\title{Percepciones meritocráticas: percepciones de esfuerzo y talento
en la distribución de recompensas sus factores asociados}
\usepackage{etoolbox}
\makeatletter
\providecommand{\subtitle}[1]{% add subtitle to \maketitle
  \apptocmd{\@title}{\par {\large #1 \par}}{}{}
}
\makeatother
\subtitle{Seminario de Grado 1 2026}
\author{Bastián Javier Martin Gerdes}
\date{26 de noviembre de 2026}
\begin{document}
\frontmatter
\maketitle

% --- title-pdf.tex personalizado para portada formal ---
\makeatletter
\providecommand{\subtitle}[1]{\gdef\@subtitle{#1}}
\providecommand{\@subtitle}{}
\providecommand{\frontmattercontext}{}
\providecommand{\advisorname}{}
\providecommand{\advisorlabel}{Profesor guia:}
\providecommand{\frontmatterlocation}{}
\InputIfFileExists{includes/cover-config.tex}{}{}

\newcommand{\PrintTitle}{%
  {\sffamily\bfseries\fontsize{24pt}{28pt}\selectfont \@title\par}%
}
\newcommand{\PrintSubtitle}{%
  \begingroup
  \edef\temp{\detokenize{\@subtitle}}%
  \ifx\temp\empty\relax
    % sin subtitulo
  \else
    {\sffamily\bfseries\large \@subtitle\par}%
  \fi
  \endgroup
}
\newcommand{\PrintAuthor}{%
  {\Large\bfseries \@author\par}%
}
\newcommand{\PrintDate}{%
  {\small \@date\par}%
}
\makeatother

\begin{titlepage}
\thispagestyle{empty}
\begin{center}
\vspace*{10mm}

% Logo o imagen institucional
\includegraphics[width=0.35\textwidth]{assets/cover.png}\par
\vspace{12mm}

% Titulo y subtitulo
\PrintTitle
\vspace{6mm}
\PrintSubtitle

\vspace{22mm}
\begingroup
\edef\temp{\detokenize{\frontmattercontext}}%
\ifx\temp\empty\relax
  % sin contexto adicional
\else
  {\normalsize \frontmattercontext\par}
\fi
\endgroup

\vspace{18mm}
\PrintAuthor

\vspace{16mm}
\begingroup
\edef\temp{\detokenize{\advisorname}}%
\ifx\temp\empty\relax
  % sin tutor
\else
  \rule{0.45\textwidth}{0.4pt}\par
  {\small \advisorlabel\ \advisorname\par}
\fi
\endgroup

\vfill
\begingroup
\edef\temp{\detokenize{\frontmatterlocation}}%
\ifx\temp\empty\relax
  % sin ubicacion
\else
  {\small \frontmatterlocation\par}
\fi
\endgroup
\PrintDate

\end{center}
\end{titlepage}

% Preliminares (numeros romanos) y estilo simple
\frontmatter
\pagestyle{scrheadings}

\setcounter{tocdepth}{2} % (o 1/3 segun prefieras)
\tableofcontents
\listoffigures
\listoftables


\mainmatter
\bookmarksetup{startatroot}

\chapter{index}\label{index}

\small\textbf{Palabras clave:} Sociología, desigualdades, educación,
Chile \normalsize

\bookmarksetup{startatroot}

\chapter*{Abstract}\label{abstract}
\addcontentsline{toc}{chapter}{Abstract}

\markboth{Abstract}{Abstract}

Resumen en español (200-300). Expón propósito, método, resultados clave
y conclusiones. Puede incluir citas como y se actualizarán
automáticamente.

\bookmarksetup{startatroot}

\chapter{Informe de Avance}\label{informe-de-avance}

\section{Problema de investigación}\label{problema-de-investigaciuxf3n}

La meritocracia es el principio de justicia más extendido en las
sociedades contemporáneas, entregando una base moral para asignar
recursos y recompensas escasas de manera proporcional al desempeño
individual, dejando de lado el origen social
(\citeproc{ref-mijs_unfulfillable_2016}{Mijs, 2016};
\citeproc{ref-mulligan_meritocracy_2023a}{Mulligan, 2023}). Además,
cumple la función de legitimar la desigualdad, ya que al reducir el
éxito a factores individuales deja de lado los factores estructurales y
presenta las brechas sociales como merecidas
(\citeproc{ref-accominotti_making_2026}{Accominotti \& Sauder, 2026};
\citeproc{ref-castillo_merit_perc_pref_2019}{Castillo et~al., 2019};
\citeproc{ref-lierse_des_redi_2022}{Lierse et~al., 2022, p. 373}). Desde
la sociología, esta legitimación se ha estudiado a través de las
percepciones meritocráticas, es decir, lo que las personas creen sobre
si la sociedad efectivamente recompensa el esfuerzo y el talento
(\citeproc{ref-castillo_merit_perc_pref_2019}{Castillo et~al., 2019}).
Estas percepciones varían según la posición que las personas ocupan y la
trayectoria que han recorrido, en particular según su estatus social
subjetivo, su nivel educacional y su experiencia de movilidad social
(\citeproc{ref-iturra-sanhueza_social_2026}{Iturra-Sanhueza et~al.,
2026}; \citeproc{ref-mijs_belief_2022}{Mijs et~al., 2022}). Sin embargo,
los factores asociados a percibir que se recompensa el esfuerzo podrían
no ser los mismos que los asociados a percibir que se recompensa el
talento.

Esta posibilidad se apoya en que ambos criterios no tienen la misma base
de merecimiento. La fórmula de Young (\citeproc{ref-young_1964}{1964}),
que define el mérito como la suma de inteligencia y esfuerzo, contiene
dos criterios que remiten a principios de merecimiento distintos
(\citeproc{ref-clavero_2023}{Clavero García, 2023};
\citeproc{ref-sardoc_rise_2024}{Sardoč, 2024}). El esfuerzo depende en
mayor medida de la voluntad y la motivación de las personas
(\citeproc{ref-franetovic_concepciones_2017a}{Franetovic, 2017, p. 5}),
mientras que el talento, que suele entenderse como un conjunto de
capacidades innatas distribuidas de manera azarosa
(\citeproc{ref-harden_2021}{Harden, 2021};
\citeproc{ref-wright_preguntas_2010}{Wright, 2010, pp. 50-51}), escapa
en mayor medida a su control, por lo que la desigualdad basada en el
talento tiende a considerarse menos legítima
(\citeproc{ref-castillo_perceptions_2025}{Castillo et~al., 2025}). Esta
diferencia también cambia el sentido de la desventaja, ya que atribuir
una posición baja a la falta de esfuerzo la convierte en una
responsabilidad individual, mientras que atribuirla a la falta de
talento la presenta como una carencia natural de la persona
(\citeproc{ref-sandel_tirania_2020}{Sandel, 2020};
\citeproc{ref-young_1964}{Young, 1964}).

Esta distinción no es solo conceptual, ya que personas de diferente
estatus social tienden a explicar, justificar y dar sentido a sus
posiciones en la estructura social basada en el mérito. En Chile, las
élites económicas explican el acceso a las posiciones superiores más por
el talento que por el esfuerzo, mientras que en las clases medias y
bajas, y en las universidades de masas, predomina el esfuerzo
(\citeproc{ref-atria_et_al_elites_meritocracia_2020}{Atria et~al.,
2020}; \citeproc{ref-quaresma_merit_u_elites_2024}{Quaresma, Madrid,
et~al., 2024}). Quienes han experimentado movilidad ascendente tienden a
destacar su esfuerzo (\citeproc{ref-jin_meritocracy_2020}{Jin \& Ball,
2020}; \citeproc{ref-mijs_belief_2022}{Mijs et~al., 2022}), aunque
también pueden reivindicar su talento para distinguirse de sus pares
privilegiados (\citeproc{ref-renteria_telling_2026}{Rentería, 2026}), y
el criterio que privilegian las élites varía según el país y las
instituciones en las que se forman
(\citeproc{ref-friedman_meaning_2024}{Friedman et~al., 2024};
\citeproc{ref-gogescu_meritocracy_2025}{Gogescu, 2025};
\citeproc{ref-warikoo_legitimating_2014}{Warikoo \& Fuhr, 2014}).

Hasta ahora, sin embargo, se han analizado las percepciones
meritocráticas predominantemente bajo la fórmula de Young
(\citeproc{ref-young_1964}{1964}) y algunos estudios las miden solo a
través del esfuerzo (\citeproc{ref-ehmes_paradox_2026}{Ehmes \& Gangl,
2026}; \citeproc{ref-mijs_paradox_2021}{Mijs, 2021};
\citeproc{ref-wiesner_effort_2026}{Wiesner \& Sachweh, 2026}) y otros
suman esfuerzo y talento en un mismo indicador
(\citeproc{ref-castillo_et_al_multidimensional_merit_beliefs_2023}{Castillo
et~al., 2023};
\citeproc{ref-iturra-sanhueza_social_2026}{Iturra-Sanhueza et~al.,
2026}), pero rara vez se analiza el talento por separado\footnote{La
  encuesta ISSP (International Social Survey Programme) incluye en su
  módulo de desigualdad social la importancia de las habilidades
  naturales para salir adelante: ``Get ahead: How important is having
  natural ability?'' (\citeproc{ref-issp_data}{ISSP Research Group,
  2002}). ELSOC (Estudio Longitudinal Social de Chile) incluye el ítem
  ``En Chile, las personas son recompensadas por su inteligencia y
  habilidad'' (\citeproc{ref-elsoc_data}{Reproducible Research, Centre
  for Social Conflict and Cohesion Studies COES, 2024}). La encuesta
  CFPS (China Family Panel Studies) pregunta ``The most important factor
  affecting one's future success is his/her talent'' y, en olas
  enfocadas en el rendimiento educativo, ``In your opinion, talent
  accounts for \_\_\_\_\_\_\% in achieving a good academic record''
  (\citeproc{ref-cfps_data}{Institute of Social Science Survey, Peking
  University, 2015}).}. Así, los factores asociados a la percepción
meritocrática se han estimado sobre medidas que no distinguen entre
ambos criterios, y no es posible saber si los resultados valen para los
dos o solo para uno de ellos.

Hay indicios de que ambos criterios se comportan de manera distinta. En
Chile, el esfuerzo se percibe como menos recompensado que la
inteligencia y las habilidades
(\citeproc{ref-castillo_merit_perc_pref_2019}{Castillo et~al., 2019, p.
6}; \citeproc{ref-iturra-sanhueza_social_2026}{Iturra-Sanhueza et~al.,
2026, p. 6}), y cuando ambas percepciones se estiman por separado, su
asociación con el estatus subjetivo y la educación difiere en magnitud,
aunque esas diferencias no han sido puestas a prueba formalmente
(\citeproc{ref-iturra-sanhueza_social_2026}{Iturra-Sanhueza et~al.,
2026}, Tabla S3). En este escenario, el presente estudio busca
establecer en qué medida los factores asociados a la percepción de que
en Chile se recompensa el esfuerzo difieren de aquellos asociados a la
percepción de que se recompensa el talento.

Para ello, se consideran tres factores asociados a la posición y a la
trayectoria social de las personas. El primero es el estatus social
subjetivo, ya que quienes se perciben en posiciones más altas tienden a
percibir que la sociedad premia el mérito, lo que legitima su propia
posición (\citeproc{ref-castillo_merit_perc_pref_2019}{Castillo et~al.,
2019}; \citeproc{ref-iturra-sanhueza_social_2026}{Iturra-Sanhueza
et~al., 2026}). El segundo es el nivel educacional, ya que la educación
se asocia con una visión más crítica de cómo se distribuyen las
recompensas, pero al mismo tiempo certifica los logros académicos a
través del mérito individual
(\citeproc{ref-castillo_merit_perc_pref_2019}{Castillo et~al., 2019};
\citeproc{ref-iturra-sanhueza_social_2026}{Iturra-Sanhueza et~al.,
2026}). El tercero es la movilidad social subjetiva, ya que quienes
perciben haber ascendido tienden a interpretar la sociedad a partir de
su propia trayectoria y a percibirla como más meritocrática
(\citeproc{ref-mijs_belief_2022}{Mijs et~al., 2022};
\citeproc{ref-traini_et_al_mov_social_percepciones_meritocracia_2025}{Traini
et~al., 2025}) y por último el nivel socioeconómico, donde quienes son
objetivamente más aventajados y elites buscan justificar su posición en
torno al mérito individual y no ventajas estructurales.

\section{Objeto de estudio}\label{objeto-de-estudio}

La investigación sobre meritocracia distingue entre las preferencias,
que son creencias normativas sobre cómo debería funcionar la sociedad, y
las percepciones, que son descripciones de cómo funciona y que pueden
diferir de lo que efectivamente sucede en ella
(\citeproc{ref-castillo_percepcion_2012}{Castillo et~al., 2012};
\citeproc{ref-castillo_merit_perc_pref_2019}{Castillo et~al., 2019}).
Percibir que una sociedad premia el esfuerzo o el talento es, por lo
tanto, distinto de preferir que lo haga. Según el teorema de Thomas, si
las personas definen una situación como real, esta es real en sus
consecuencias (\citeproc{ref-merton_selffulfilling_1948}{Merton, 1948}).
Así, percibir que la sociedad es meritocrática se asocia con considerar
más aceptables las desigualdades (\citeproc{ref-dubet_nuevo_2023}{Dubet,
2023, pp. 32-33}) y con un menor apoyo a la redistribución
(\citeproc{ref-sonhing_merit_2011}{Son Hing et~al., 2011};
\citeproc{ref-tejero-peregrina_perceived_2025}{Tejero-Peregrina et~al.,
2025}). Además, estas consecuencias difieren según el criterio. En
Chile, percibir que se recompensa el esfuerzo se asocia con legitimar
que el acceso a la salud, la educación y las pensiones dependa del
ingreso, mientras que percibir que se recompensa el talento no se
asocia, o se asocia negativamente, con esa legitimación
(\citeproc{ref-castillo_perceptions_2025}{Castillo et~al., 2025},
\citeproc{ref-castillo_justification_desigualdades_pensiones_clases_2026}{2026}).

Así, esta investigación tiene como objeto de estudio las percepciones
sobre la meritocracia, diferenciando entre la percepción de que se
recompensa el esfuerzo y la percepción de que se recompensa el talento.

\section{Pregunta de investigación}\label{pregunta-de-investigaciuxf3n}

¿En qué medida el estatus social subjetivo, el nivel educacional, el
nivel socioeconómico y la experiencia de movilidad social se asocian con
la percepción de que en Chile se recompensa el esfuerzo y con la
percepción de que se recompensa el talento, y en qué medida difieren
estas asociaciones entre ambas percepciones?

\section{Marco conceptual}\label{marco-conceptual}

\subsection{Meritocracia: concepto y
debate}\label{meritocracia-concepto-y-debate}

El concepto de meritocracia se populariza con Young
(\citeproc{ref-young_1964}{1964}) como una sátira a una sociedad
gobernada por élites, sustentada en la fórmula mérito = inteligencia +
esfuerzo (M = I + E). La idea supone la igualdad de oportunidades, en la
cual los individuos, a través del esfuerzo y el talento, pueden lograr
movilidad social y éxito en ciertas esferas sociales, dejando de lado
factores no meritocráticos como el género, la clase social o el origen
familiar (\citeproc{ref-wooldridge_aristocracy_2021}{Wooldridge, 2021}).
Aun con la carga negativa que le atribuía
(\citeproc{ref-young_1964}{Young, 1964}; también en
\citeproc{ref-young_2001}{Young, 2001}), el concepto ha sido adoptado
como ideal normativo en distintas sociedades para distribuir recompensas
y recursos, seleccionar personas y legitimar las diferencias en la
educación, el mundo empresarial y la política
(\citeproc{ref-mijs_unfulfillable_2016}{Mijs, 2016}).

En la sociología el mérito es un tipo de merecimiento propio de la
justicia distributiva, que depende de las normas, valores y objetivos de
prácticas sociales particulares (\citeproc{ref-pavel_merit_2024}{Pavel,
2024}).El concepto en sí no es sinónimo de igualdad de oportunidades,
pero se utiliza de manera amplia en el debate político, puesto que es el
principio más cercano a la idea de carreras abiertas al talento
(\citeproc{ref-sardoc_rise_2024}{Sardoč, 2024}). De esta manera, el
mérito ofrece una base moral para la asignación de bienes escasos, ya
que se recompensa lo que se hace, y quien se desempeña mejor merece una
recompensa mayor (\citeproc{ref-mijs_unfulfillable_2016}{Mijs, 2016}).
Según Mijs (\citeproc{ref-mijs_unfulfillable_2016}{2016}), la
meritocracia se compone de tres elementos fundamentales: las carreras
abiertas al talento, que asignan los puestos a quienes demuestran
competencia y no a través del nepotismo o el soborno; las oportunidades
educativas adaptadas a la capacidad natural, de modo que los procesos de
prueba y selección del sistema educativo desarrollen a quienes tienen
potencial y no se desperdicien talentos; y el logro como base de la
desigualdad social, que reemplaza a la adscripción para que la sociedad
funcione de manera más eficiente. El atractivo de la meritocracia para
los responsables políticos reside en esta combinación que permitiría
legitimar las diferencias, estimular el esfuerzo y optimizar la
asignación de recompensas.

En la misma línea, Sen (\citeproc{ref-sen_merit_justice_2000}{2000})
entiende que el término está subdefinido y que no hay mucha claridad
sobre el significado de meritocracia, al mismo tiempo, menciona que
premiar el mérito significa premiar los criterios sociales normativos de
éxito, según la concepción de buena sociedad que tenga esta misma lo que
permite un procedimiento legítimo para asignar recursos y recompensas
(\textbf{o simplemente decir lo que dice Rawls que está basado en
factores arbitrarios como estandares culturales cambiantes(e innatos)}).
Pero, para poder determinar lo que es meritorio y lo que no, se necesita
un grado de acuerdo sobre principios de justicia. Esto concuerda con la
discusión reciente de Clavero García (\citeproc{ref-clavero_2023}{2023,
p. 1180}), quien define el mérito como ``the value of a free, effortful,
good-oriented action'' {[}el valor de una acción libre, esforzada y
orientada al bien{]}, lo que lo vuelve contextual y diferente según la
sociedad. En su formulación clásica, el mérito combina capacidad o
talento con esfuerzo, y este último incluye la dedicación y el trabajo
sostenido hacia un fin (\citeproc{ref-clavero_2023}{Clavero García,
2023}; \citeproc{ref-franetovic_concepciones_2017a}{Franetovic, 2017}).
Bajo esta definición, el talento solo sería meritorio en la medida en
que se ejerce mediante una acción libre y esforzada, y no por sí mismo,
lo que tensiona la fórmula de Young (\citeproc{ref-young_1964}{1964}),
que lo incluye como componente directo del mérito.

Sobre el ideal, el debate ha estado en enfocado en meritócratas como
Wooldridge (\citeproc{ref-wooldridge_aristocracy_2021}{2021}) quien
considera que la meritocracia actual está incompleta y que debe seguir
desarrollándose para superar sus problemas, ya que permite juzgar a las
personas por sus méritos en lugar de las ventajas heredadas, el
nepotismo o el clientelismo, algo que han valorado y exigido los
sectores históricamente desfavorecidos. En cambio, para Sandel
(\citeproc{ref-sandel_tirania_2020}{2020}) el problema no radica solo en
que la igualdad de oportunidades no se haya alcanzado, sino en el
principio mismo, puesto que al hacer creer que las posiciones sociales
son el resultado del esfuerzo y el talento individual genera soberbia en
quienes triunfan y humillación en quienes quedan rezagados, y convierte
los títulos en una medida del valor de las personas.

El debate normativo, sin embargo, no permite responder si la sociedad
funciona de manera meritoria. Esa pregunta ha sido abordada por la
sociología desde la teoría funcionalista de K. Davis \& Moore
(\citeproc{ref-davis_est_1970}{1970}), según la cual la desigualdad
opera como un esquema de motivación donde el éxito y las recompensas
dependen de dos componentes del mérito, la capacidad o talento
individual y el esfuerzo o capacitación, por lo que sería justa e
indispensable para que los más aptos asuman los roles de mayor
responsabilidad. Además, en las sociedades modernas las características
adquiridas serían más relevantes que las adscritas
(\citeproc{ref-parsons_class_1953}{Parsons, 1953}) Pero, que las
sociedades funcionen efectivamente de manera meritocrática ha mantenido
críticas, ya que, se puede apreciar cómo organizaciones educativas no
responden a la necesidad funcional o al mérito técnico
(\citeproc{ref-hauhart_davis_2003}{Hauhart, 2003}), sino al poder y
grupos de interés. De esta manera la promesa de la meritocracia basada
en la educación no permite la movilidad social de las clases bajas y
medias, según la cual las posiciones de clase dependerían cada vez más
de las credenciales educativas y menos del origen social, mostrando sólo
pequeños cambios (\citeproc{ref-goldthorpe_myth_2003}{Goldthorpe,
2003}).

Más allá de si la meritocracia se cumple, la sociología crítica ha
planteado que opera como un mecanismo de legitimación. Para Bourdieu \&
Passeron (\citeproc{ref-bourdieu_reproduccion_1998}{1998}), el sistema
educativo ejerce una violencia simbólica al presentar como méritos
individuales diferencias que dependen del origen social, con lo que las
desigualdades aparecen como diferencias naturales entre las personas.

De este modo se presentarían las desigualdades como diferencias
naturales entre los individuos y no de sus distintas condiciones
sociales. En esta línea Accominotti \& Sauder
(\citeproc{ref-accominotti_making_2026}{2026}) plantean que la
meritocracia no solo ordena a las personas en posiciones desiguales,
sino que además construye órdenes de mérito, es decir, jerarquías de
estatus basadas en evaluaciones del mérito, el logro o el desempeño, que
establecen relaciones de superioridad, igualdad o inferioridad entre las
personas. Estos órdenes tienen efectos propios, ya que moralizan la
desigualdad al presentar las diferencias de ventajas como diferencias de
mérito individual y, en línea con Sandel
(\citeproc{ref-sandel_tirania_2020}{2020}), pueden erosionar la cohesión
social y la solidaridad distributiva al reintroducir jerarquías morales
sobre el valor de las personas. Además, existen como esquemas culturales
más o menos compartidos, de modo que lo que las personas perciben como
merecido forma parte de ellos.

En el estudio de las creencias meritocráticas se distinguen dos
enfoques, por un lado, las preferencias, que son creencias normativas
sobre cómo debería funcionar la sociedad, y por otro, las percepciones,
que son creencias descriptivas sobre cómo funciona efectivamente
(\citeproc{ref-castillo_merit_perc_pref_2019}{Castillo et~al., 2019}).
Como menciona Castillo et~al.
(\citeproc{ref-castillo_merit_perc_pref_2019}{2019}), percibir que la
meritocracia funciona no implica considerarla el mejor sistema ni
preferirla. Además, las percepciones cumplen una función legitimadora,
puesto que quienes perciben que la meritocracia funciona también
perciben menos desigualdad económica
(\citeproc{ref-castillo_merit_perc_pref_2019}{Castillo et~al., 2019, p.
11}). Estas percepciones no son descripciones neutras de la sociedad,
sino inferencias que las personas construyen a partir de su propia
trayectoria y de su entorno social, que permiten entender cómo es la
meritocracia en la sociedad (\citeproc{ref-dawtry_why_2015}{Dawtry
et~al., 2015}; \citeproc{ref-mijs_inequality_2018}{J. Mijs, 2018}).
Además, las personas reconocen factores meritocráticos y no
meritocráticos, como el origen familiar o los contactos, lo que se ha
denominado conciencia dual
(\citeproc{ref-reynolds_perceptions_2014}{Reynolds \& Xian, 2014};
\citeproc{ref-wiesner_effort_2026}{Wiesner \& Sachweh, 2026}). Por ello,
las percepciones meritocráticas y no meritocráticas se miden como
dimensiones distintas
(\citeproc{ref-castillo_et_al_multidimensional_merit_beliefs_2023}{Castillo
et~al., 2023}), y la forma en que se combinan varía entre países
(\citeproc{ref-zhu_meritocratic_2025}{Zhu, 2025}).

\subsubsection{Esfuerzo y talento}\label{esfuerzo-y-talento}

Como se señaló, tanto la fórmula de Young
(\citeproc{ref-young_1964}{1964}) como la teoría de K. Davis \& Moore
(\citeproc{ref-davis_est_1970}{1970}) consideran el talento y el
esfuerzo como los dos componentes del mérito, que son también las
explicaciones más comunes del éxito y del fracaso
(\citeproc{ref-weiner_attributional_1985}{Weiner, 1985, p. 549}). Sin
embargo, ambos no se sostienen en la misma base de merecimiento, ya que
el esfuerzo se percibe como algo que depende de las personas, mientras
que el talento ocupa una posición ambigua entre el logro y la herencia.
Por ello, percibir que la sociedad recompensa el esfuerzo implica creer
que las recompensas dependen de lo que las personas hacen, mientras que
percibir que recompensa el talento implica creer que dependen de lo que
las personas tienen.

Además, las personas pueden percibir el esfuerzo y el talento de
distintas maneras, ya sea según su importancia para salir adelante, como
explicación de la pobreza y la riqueza, o como explicación de la propia
posición. Esta investigación se centra en la percepción de que la
sociedad recompensa a quienes más se esfuerzan y a quienes poseen más
talento (\citeproc{ref-castillo_merit_perc_pref_2019}{Castillo et~al.,
2019}; \citeproc{ref-castillo_perceptions_2025}{Castillo et~al., 2025}).
A continuación se analizan ambos criterios por separado.

\paragraph{Esfuerzo}\label{esfuerzo}

La ética del esfuerzo ha sido un tema recurrente en la sociología, en
específico en la distribución de recursos. Desde Weber
(\citeproc{ref-weber_etica_2001}{2001}) se ha analizado el esfuerzo en
relación con la religión, en torno a la secularización y la ética
protestante del trabajo, donde se tensiona la idea de una omnipotencia
divina que predetermina el destino, como con el sistema católico de
salvación mediante obras e indulgencias. Así, el trabajo duro y el
ascetismo pasan a ser una señal de salvación, permitiendo tener
respuestas a la incertidumbre psicológica de la predestinación. Con la
secularización de los valores religiosos, el orden meritocrático
moraliza el éxito a través del esfuerzo personal, trasladando el énfasis
en el libre albedrío y agencia de las personas
(\citeproc{ref-sandel_tirania_2020}{Sandel, 2020, pp. 41-60}). De este
modo, el éxito se vuelve una señal de esfuerzo, y el esfuerzo, el
fundamento para considerar merecidas las recompensas.

Por otra parte, Merton (\citeproc{ref-merton_social_1938}{1938}) planteó
que la sociedad estadounidense sostenía una ideología de clases
abiertas, según la cual el éxito era una meta deseable y alcanzable para
todos, y que esta ideología persistía porque contribuía a mantener el
orden existente\footnote{Se retoma solo la descripción que hace Merton
  de esta ideología y no su teoría de la desviación, que ha sido
  criticada por suponer que todos los grupos comparten las mismas metas
  culturales.}. En una línea parecida, Kluegel \& Smith
(\citeproc{ref-kluegel_beliefs_1981}{1981}) denominaron ideología
dominante al conjunto de creencias según las cuales las oportunidades
están ampliamente disponibles y cada persona es responsable de su
posición, por lo que la desigualdad resultante es justa, y el esfuerzo
individual aparece como el camino legítimo hacia el éxito, y su falta,
como la explicación del fracaso.

Las teorías que siguen caracterizan qué implica explicar un resultado
por el esfuerzo o por el talento, y no miden directamente la percepción
de que la sociedad los recompensa. Sin embargo, permiten entender por
qué ambas percepciones podrían asociarse de forma distinta con la
posición de las personas. Desde la teoría de la atribución, Weiner
(\citeproc{ref-weiner_attributional_1985}{1985}) plantea que las
personas buscan las causas de sus éxitos y fracasos para comprender su
entorno y orientar sus acciones futuras. El éxito suele atribuirse a una
alta capacidad y al trabajo duro, y el fracaso, a una baja capacidad y a
la falta de esfuerzo. El autor clasifica estas causas según su locus, su
estabilidad y su controlabilidad. El esfuerzo y la capacidad son causas
internas, pero solo el esfuerzo se considera controlable, y es esta
controlabilidad la que lleva a juzgar a una persona como responsable de
sus resultados, ya sea al evaluar a otros o a uno mismo
(\citeproc{ref-davis_learning_2017}{M. L. Davis et~al., 2017};
\citeproc{ref-schrader_school_2015}{Schrader \& Helmke, 2015}). Lo que
distingue al esfuerzo de otros criterios de mérito es, entonces, que se
percibe bajo el control de las personas, ya que depende en mayor medida
de su voluntad y motivación
(\citeproc{ref-franetovic_concepciones_2017a}{Franetovic, 2017, p. 5}).
Este control es uno de los criterios con que las personas juzgan lo que
otros merecen (\citeproc{ref-meuleman_desert_carin_2020}{Meuleman
et~al., 2020}), y, en Estados Unidos, el esfuerzo es percibido como el
factor que más depende de las personas y se juzga como el criterio más
justo para determinar el ingreso
(\citeproc{ref-bartling_talent_2026}{Bartling et~al., 2026, p. 6}).

Percibir el esfuerzo como algo controlable tendría consecuencias por
atribuirlo a factores propios o externos. Quien se esfuerza merece ser
recompensado por su dedicación, pero también se usa para explicar los
fracasos, con componentes de culpa o desesperanza, porque supone que la
situación estaba bajo el control de la persona
(\citeproc{ref-lampert_meritocratic_2013}{Lampert, 2013};
\citeproc{ref-spruyt_talent_esfuerzo_social_background_2015}{Spruyt,
2015, p. 96}). Así, un fracaso atribuido a la falta de esfuerzo genera
enojo o rechazo en otros y culpa en uno mismo, mientras que uno
atribuido a la falta de capacidad genera lástima y vergüenza
(\citeproc{ref-weiner_attributional_1985}{Weiner, 1985, p. 561}).

Esta distinción entre culpa y destino también organiza las explicaciones
de la pobreza, que se clasifican según si la causa está en el individuo
o en la sociedad, y si remite a la acción de alguien que puede ser
culpado o a una cuestión del destino
(\citeproc{ref-oorschot_blame_2000}{Oorschot \& Halman, 2000, pp. 5-7}).
La explicación de culpa individual atribuye la pobreza a fallas morales
o conductuales de los propios pobres, como pereza, falta de ahorro o
malas costumbres, mientras que la de destino individual la atribuye a la
mala suerte, la voluntad divina o infortunios personales que están fuera
del control propio (\citeproc{ref-oorschot_blame_2000}{Oorschot \&
Halman, 2000, p. 7}). De este modo, una posición baja atribuida a la
falta de esfuerzo se convierte en una responsabilidad individual
(\citeproc{ref-sandel_tirania_2020}{Sandel, 2020}).

El talento, en cambio, ocupa un lugar más ambiguo. En los estudios que
revisan Oorschot \& Halman (\citeproc{ref-oorschot_blame_2000}{2000, p.
6}), la falta de habilidad y talento se agrupa con la falta de esfuerzo
entre las explicaciones individuales de la pobreza. Sin embargo, el
talento se entiende como una cualidad estable cuya presencia escapa al
control de las personas, por lo que explicar un resultado por el talento
resulta más fatalista que culpabilizador
(\citeproc{ref-spruyt_talent_esfuerzo_social_background_2015}{Spruyt,
2015, pp. 96-97}).

Además, los límites del esfuerzo están en relación con la estructura
social, por barreras materiales o culturales de origen. Así, Bourdieu
(\citeproc{ref-bourdieu_distincion_2006}{2006}) describe cómo la pequeña
burguesía se caracteriza por una buena voluntad cultural, es decir, por
el esfuerzo de apropiarse de la cultura legítima, mientras que las
clases dominantes valoran la soltura y la naturalidad, que hacen
aparecer como innato lo que fue adquirido en la familia. Del mismo modo,
en contextos escolares se valora el esfuerzo, pero lograr buenos
resultados sin esfuerzo se celebra como señal de una inteligencia
auténtica (\citeproc{ref-jackson_smart_2015}{Jackson \& Nyström, 2015,
pp. 393-394}). A su vez, el esfuerzo que realizan las familias de clase
trabajadora para compensar su menor capital cultural no siempre es
recompensado (\citeproc{ref-grau-grau_ilusion_2022}{Grau-Grau, 2022, p.
4}), mientras que quienes heredaron mayor capital cultural declaran
haberse esforzado más que el resto para obtener su posición
(\citeproc{ref-grau-grau_ilusion_2022}{Grau-Grau, 2022, pp. 15-16}). En
el mercado laboral chileno, además, el esfuerzo excede al mérito, ya que
incluye la disponibilidad, la personalidad y las relaciones personales,
en un contexto donde hay un desconocimiento generalizado sobre los
criterios de evaluación y asignación de recompensas
(\citeproc{ref-sir_esfuerzo_2017}{Sir, 2017, pp. 313, 319-320}).

En Chile, las trayectorias están marcadas por relatos de esfuerzo, que
operan como principio justificador, como fuente de autovaloración y como
una expectativa concreta desde la cual se evalúa la justicia de la
sociedad y el reconocimiento
(\citeproc{ref-araujo_y_martuccelli_soc-chilena_2012}{Araujo \&
Martuccelli, 2012}). En el mundo del trabajo, el esfuerzo se ha
convertido en el núcleo ético de un mercado laboral individualizado, al
punto de que ``el esfuerzo y no el mérito es lo que reclama más
compensación'' (\citeproc{ref-sir_esfuerzo_2017}{Sir, 2017, p. 311}), y
en un contexto de desprotección laboral se vuelve un requisito para ser
reconocido (\citeproc{ref-sir_esfuerzo_2017}{Sir, 2017, p. 320}).

Percibir que la sociedad recompensa el esfuerzo implica, entonces, creer
que las recompensas siguen a lo que las personas hacen, una percepción
que se alinea con los relatos sobre la responsabilidad individual y el
merecimiento moral (\citeproc{ref-castillo_perceptions_2025}{Castillo
et~al., 2025}). En Chile, esta es la percepción que más se asocia con la
legitimación de la desigualdad, ya que percibir que se recompensa el
esfuerzo se asocia con aceptar que el acceso a la salud, la educación y
las pensiones dependa del ingreso
(\citeproc{ref-castillo_perceptions_2025}{Castillo et~al., 2025},
\citeproc{ref-castillo_justification_desigualdades_pensiones_clases_2026}{2026}).
Una explicación posible es que el esfuerzo se ve como un rasgo
controlable, mientras que el talento se percibe como innato y menos
sujeto al control personal
(\citeproc{ref-castillo_perceptions_2025}{Castillo et~al., 2025}).

\paragraph{Talento}\label{talento}

En esta investigación, el talento se entiende como una noción de sentido
común sobre las capacidades que las personas tienen, sin pronunciarse
sobre si estas son innatas o adquiridas. Por ello, interesa qué
reconocen las personas como talento cuando evalúan si la sociedad lo
recompensa.

Desde la sociología, el talento ha sido entendido como algo que no es
independiente de la posición social. Bourdieu \& Passeron
(\citeproc{ref-bourdieu_2009_herederos}{2009}) describen una ideología
del don, según la cual el éxito escolar de los estudiantes de origen
privilegiado se percibe como resultado de dotes naturales y no de una
herencia cultural transmitida por la familia. En su teoría de la
reproducción, Bourdieu \& Passeron
(\citeproc{ref-bourdieu_reproduccion_1998}{1998}) sostiene que las
desigualdades educativas no nacen del mérito y el talento individual,
sino del funcionamiento de campos que premian las prácticas de quienes
están familiarizados con la cultura dominante
(\citeproc{ref-grau-grau_ilusion_2022}{Grau-Grau, 2022, p. 3}). En este
marco, el capital cultural incorporado se adquiere a través de la
socialización en la familia de origen y se expresa en disposiciones, es
decir, en formas de hablar, juzgar, percibir y actuar que forman parte
del habitus de cada persona
(\citeproc{ref-bourdieu_distincion_2006}{Bourdieu, 2006};
\citeproc{ref-grau-grau_ilusion_2022}{Grau-Grau, 2022}). El sistema
escolar, bajo una apariencia de neutralidad, valora estas disposiciones
como si fueran habilidades personales, por lo que termina premiando a
quienes ya estaban aventajados y transformando un privilegio heredado en
mérito propio (\citeproc{ref-grau-grau_ilusion_2022}{Grau-Grau, 2022}).
Desde esta perspectiva, lo que se reconoce como talento puede ser una
ventaja de origen que se naturaliza y se presenta como una cualidad
individual.

Ahora bien, esto no significa negar las diferencias entre las personas,
ya que existen diferencias genéticas entre los individuos, que están
vinculadas con resultados como el nivel educativo alcanzado y que
otorgan ventajas a quienes las heredan, aunque su distribución en la
sociedad sea producto del azar (\citeproc{ref-harden_2021}{Harden,
2021}). Mijs (\citeproc{ref-mijs_unfulfillable_2016}{2016, p. 19})
desarrolla la misma idea, mencionando que las condiciones para tener
mérito están determinadas por esta lotería de nacimiento, que incluye
los recursos familiares heredados y habilidades naturales como el
coeficiente intelectual, la condición física y el atractivo. De esta
manera, el ideal meritocrático parte de condiciones desiguales que son
ajenas a la meritocracia, lo que vulnera el propio principio
meritocrático (\citeproc{ref-mijs_unfulfillable_2016}{Mijs, 2016}).
Además, no está claro qué es una habilidad ni cómo identificar las que
requiere cada tarea, aunque en ámbitos como el arte o el deporte sí es
posible aislar habilidades específicas
(\citeproc{ref-lampert_meritocratic_2013}{Lampert, 2013, p. 53}). Por lo
tanto, no es preciso discutir si el talento existe o no, sino qué se
reconoce como talento y cómo esta percepción se usa para explicar
posiciones y recompensas sociales.

Lo que se percibe como talento, sin embargo, no siempre corresponde a
una habilidad innata, según Chambliss
(\citeproc{ref-chambliss_mundanity_1989}{1989}), en su estudio
etnográfico sobre nadadores de distintos niveles, que la excelencia no
se puede explicar por un don especial, sino más bien por la acumulación
de prácticas, hábitos y detalles cotidianos que se aprenden y se llevan
a cabo de manera constante. Según el autor, recurrir al talento como
explicación, en última instancia, mistifica la forma en que realmente
surgen los logros sobresalientes, lo que en la práctica es el resultado
del entrenamiento y de un entorno propicio. De esta manera, el talento
no solo sería la causa del éxito, sino también la forma de nombrarlo una
vez que ya ha ocurrido. En el mundo del trabajo británico, las formas de
presentarse de las clases privilegiadas, como ciertos modos de hablar,
de vestir o de desenvolverse, suelen ser reconocidas como talento, lo
que explica parte de las ventajas de quienes provienen de esos orígenes
en las ocupaciones de élite
(\citeproc{ref-friedman_class_2019b}{Friedman \& Laurison, 2019}). En
Chile, las autoridades de universidades de élite asumen que las
credenciales académicas y los puntajes de admisión reflejan el talento y
el mérito de sus estudiantes
(\citeproc{ref-quaresma_justify_2024}{Quaresma, Villalobos, et~al.,
2024, p. 9}), aun cuando los programas de élite son socialmente
homogéneos (\citeproc{ref-quaresma_justify_2024}{Quaresma, Villalobos,
et~al., 2024, p. 13}).

Las personas también recurren al talento para justificar su propia
posición, en especial quienes ocupan posiciones de élite. Entre las
élites británicas predomina una justificación basada en la aptitud
natural, la inteligencia y el buen juicio innato, mientras que las
élites danesas destacan el trabajo duro, diferencia que los autores
vinculan con el papel de las escuelas privadas de élite británicas en
cultivar la idea de un talento excepcional
(\citeproc{ref-friedman_meaning_2024}{Friedman et~al., 2024, p. 861}).
Además, entre las pocas personas de la élite británica que eran mujeres,
no blancas o de origen de clase trabajadora, se destacaba más el trabajo
duro que el talento como explicación de su éxito
(\citeproc{ref-friedman_meaning_2024}{Friedman et~al., 2024, p. 872}).
De manera similar, estudiantes de Oxford reconocen que otros tienen
menos acceso a las habilidades que se cultivan en los colegios privados,
pero atribuyen su propia admisión a una inteligencia que consideran
inherente (\citeproc{ref-warikoo_legitimating_2014}{Warikoo \& Fuhr,
2014, p. 712}). En la región, élites económicas chilenas definen el
talento como habilidades empresariales y de liderazgo, más que como
inteligencia, y consideran que está disponible para todos, ya que se
trata de habilidades que cualquiera podría identificar y reproducir
(\citeproc{ref-atria_et_al_elites_meritocracia_2020}{Atria et~al., 2020,
pp. 10-11}). En Lima, en cambio, profesionales que ascendieron
socialmente reivindican su talento para distinguirse de sus pares
privilegiados, estableciendo una frontera simbólica entre el privilegio
y el talento (\citeproc{ref-renteria_telling_2026}{Rentería, 2026, pp.
1342, 1354}).

En conjunto, estos relatos muestran que el criterio que se enfatiza para
justificar la propia posición depende tanto de la posición que se ocupa
como de la socialización en instituciones como las escuelas de élite o
el mundo empresarial
(\citeproc{ref-atria_et_al_elites_meritocracia_2020}{Atria et~al.,
2020}; \citeproc{ref-friedman_meaning_2024}{Friedman et~al., 2024};
\citeproc{ref-warikoo_legitimating_2014}{Warikoo \& Fuhr, 2014}). Se
trata, sin embargo, de relatos de personas que ocupan objetivamente
posiciones aventajadas, por lo que su relación con la percepción de
recompensa según el estatus subjetivo se discute en los antecedentes.
Además, como muestra el caso de las élites chilenas, lo que se entiende
por talento varía, por lo que al evaluar si la sociedad lo recompensa
las personas pueden estar pensando en cualidades distintas, lo que
constituye una limitación para interpretar esta percepción.

Percibir que la sociedad recompensa el talento implica, entonces, creer
que las recompensas siguen a lo que las personas tienen. Por ejemplo, en
Estados Unidos, las personas sitúan el talento en algún punto entre el
esfuerzo y la suerte, ya que lo perciben como menos controlable que el
esfuerzo, pero mucho más que la suerte, y consideran que los ingresos
derivados del talento son casi tan justos como los que resultan del
esfuerzo (\citeproc{ref-bartling_talent_2026}{Bartling et~al., 2026, p.
6}). Según Bartling et~al. (\citeproc{ref-bartling_talent_2026}{2026}),
esto se debe a que, cuando una persona hace uso de su talento, se la
considera dueña de sus resultados, aunque el talento no dependa de ella.
Sin embargo, atribuir el éxito a un talento innato puede presentar la
desigualdad como inevitable, como resultado de diferencias naturales de
capacidad (\citeproc{ref-friedman_meaning_2024}{Friedman et~al., 2024}).
En Chile, a diferencia de lo que ocurre con el esfuerzo, percibir que se
recompensa el talento no se asocia, o se asocia negativamente, con
legitimar que el acceso a la salud, la educación y las pensiones dependa
del ingreso (\citeproc{ref-castillo_perceptions_2025}{Castillo et~al.,
2025},
\citeproc{ref-castillo_justification_desigualdades_pensiones_clases_2026}{2026}).
A su vez, la creencia de que la sociedad recompensa el talento tiene
consecuencias para quienes se encuentran en las posiciones más bajas, ya
que atribuir una posición baja a la falta de talento la presenta como un
defecto natural del individuo, sobre el cual tiene poco control
(\citeproc{ref-sandel_tirania_2020}{Sandel, 2020};
\citeproc{ref-young_1964}{Young, 1964}).

\subsection{Esfuerzo y Talento sus factores
asociados}\label{esfuerzo-y-talento-sus-factores-asociados}

\subsubsection{Estatus social subjetivo}\label{estatus-social-subjetivo}

A diferencia de la clasificación objetiva de clase, el estatus social
subjetivo refiere a la autoubicación de los individuos en la jerarquía
social, entendida por J. A. Davis
(\citeproc{ref-davis_status_1956}{1956}) como una creencia personal que
no necesariamente coincide con su posición material real. Desde la tesis
de los grupos de referencia, este desajuste se explica porque las
personas suelen desenvolverse en entornos socioeconómicamente
homogéneos, lo que actúa como un filtro cognitivo al evaluar la
desigualdad general. Así, aunque las condiciones económicas objetivas
inciden en la autopercepción
(\citeproc{ref-iturra_status_subj_2019}{Iturra-Sanhueza \& Riffo,
2019}), su efecto es más débil que el de la comparación cotidiana con el
entorno cercano (\citeproc{ref-evans_images_1992}{Evans et~al., 1992}).
En Chile, las personas tienden a ubicarse en las posiciones medias de la
escala independientemente de su posición objetiva
(\citeproc{ref-castillo_todos_clase_media_2013}{Castillo et~al., 2013}).
Además, en contextos de alta desigualdad, quienes tienen menores
ingresos tienden a ubicarse con mayor precisión, mientras que quienes
tienen mayores ingresos tienden a subestimar su posición
(\citeproc{ref-wiesner_knowing_2026}{Wiesner \& Iturra-Sanhueza, 2026}).
Por eso, el estatus subjetivo permite captar cómo las personas se ven a
sí mismas en la jerarquía, lo cual es relevante para estudiar
percepciones, que también son evaluaciones subjetivas sobre la sociedad
y en específico cómo se distribuyen las recompensas.

La relevancia de estudiar el estatus subjetivo radica en que, a
diferencia de la posición objetiva, se asocia con mayor fuerza con la
percepción de desigualdad (\citeproc{ref-melli_uncovering_2025}{Melli \&
Azzollini, 2025}) y con los juicios sobre la justicia del propio ingreso
(\citeproc{ref-mac_baroz_merit_des_2024}{Mac-Clure et~al., 2024, p.
20}). En Chile, la evidencia cuantitativa muestra que un mayor estatus
social subjetivo se asocia con una mayor percepción de meritocracia,
incluso al controlar por el ingreso del hogar, que no resulta
significativo (\citeproc{ref-castillo_merit_perc_pref_2019}{Castillo
et~al., 2019, p. 10};
\citeproc{ref-iturra-sanhueza_social_2026}{Iturra-Sanhueza et~al.,
2026}), y que quienes se ubican en estratos subjetivos bajos evalúan su
propio ingreso como más injusto, diferencia que se acentúa cuando lo
juzgan según su esfuerzo por educarse
(\citeproc{ref-mac_baroz_merit_des_2024}{Mac-Clure et~al., 2024}). En
América Latina, un mayor estatus subjetivo también se asocia con
explicaciones más individuales de la riqueza y la pobreza
(\citeproc{ref-bucca_merit_2016}{Bucca, 2016}). No obstante, con datos
de 2016 el estatus subjetivo no se asoció significativamente con la
recompensa percibida, mientras que el estatus subjetivo de la familia de
origen se asoció positivamente y el estatus esperado para los hijos,
negativamente (\citeproc{ref-iturrasanhueza_estatus_2019}{Iturra
Sanhueza, 2019, p. 46}).

Cuando la percepción de esfuerzo y la de talento se analizan por
separado, la asociación del estatus subjetivo es positiva para la
percepción de que se recompensa el esfuerzo y para la de que se
recompensa el talento, pero de mayor magnitud en el caso del esfuerzo
(\citeproc{ref-iturra-sanhueza_social_2026}{Iturra-Sanhueza et~al.,
2026}, Tabla S3). Esto es coherente con que el esfuerzo sea el criterio
que más se asocia con la justificación de la desigualdad en Chile
(\citeproc{ref-castillo_perceptions_2025}{Castillo et~al., 2025},
\citeproc{ref-castillo_justification_desigualdades_pensiones_clases_2026}{2026}).

Las investigaciones cualitativas en Chile muestran que las élites
económicas explican el acceso a las posiciones superiores más por el
talento que por el esfuerzo
(\citeproc{ref-atria_et_al_elites_meritocracia_2020}{Atria et~al.,
2020}), y que en las universidades de élite el mérito se entiende
principalmente como talento
(\citeproc{ref-quaresma_merit_u_elites_2024}{Quaresma, Madrid, et~al.,
2024}). Sin embargo, estos estudios se refieren a personas que ocupan
objetivamente posiciones de élite, y no necesariamente a quienes se
perciben en la parte alta de la jerarquía. Dado que quienes tienen
mayores ingresos tienden a subestimar su posición
(\citeproc{ref-wiesner_knowing_2026}{Wiesner \& Iturra-Sanhueza, 2026}),
el estatus subjetivo no captura directamente a estas élites.

El estudio de Friedman et~al.
(\citeproc{ref-friedman_deflecting_2021a}{2021}) ilustra este desajuste.
Profesionales y directivos que crecieron en hogares de clase media se
identifican como clase trabajadora o como personas que ascendieron desde
abajo, a través de relatos sobre su origen que minimizan las ventajas
con las que crecieron y se apoyan en la historia de movilidad ascendente
de su familia extendida. De esta manera, aunque no han experimentado
movilidad social, presentan su trayectoria como una historia de trabajo
duro, talento excepcional y lucha meritocrática ``contra todo
pronóstico'' (\citeproc{ref-friedman_deflecting_2021a}{Friedman et~al.,
2021, p. 723}). Así, quienes ocupan posiciones aventajadas pero se
ubican más abajo de lo que indica su posición objetiva pueden justificar
su posición tanto por el esfuerzo como por el talento.

En Chile, los datos sugieren que la asociación entre el estatus
subjetivo y la percepción de que se recompensa la inteligencia y las
habilidades, ítem que se puede usar como aproximación al talento,
aumenta con el ingreso, de modo que entre quienes tienen mayores
ingresos se asemeja a la del esfuerzo, mientras que entre quienes tienen
menores ingresos la diferencia a favor del esfuerzo es mayor
(\citeproc{ref-iturra-sanhueza_social_2026}{Iturra-Sanhueza et~al.,
2026}, Tabla S3). Por ello, se espera que ambas percepciones aumenten
con el estatus subjetivo, aunque con mayor fuerza en el caso del
esfuerzo.

\begin{quote}
\(H_{1a}\): A mayor estatus social subjetivo, mayor es la percepción de
que en Chile se recompensa el esfuerzo.

\(H_{1b}\): A mayor estatus social subjetivo, mayor es la percepción de
que en Chile se recompensa el talento.

\(H_{1c}\): La asociación del estatus social subjetivo es más fuerte con
la percepción de que se recompensa el esfuerzo que con la de que se
recompensa el talento.
\end{quote}

\subsubsection{Nivel educacional}\label{nivel-educacional}

La educación ocupa un lugar central en la meritocracia, ya que es el
principal mecanismo a través del cual se seleccionan y certifican los
méritos individuales, de este modo el logro académico funciona como
indicador objetivo del mérito (Goldthorpe, 2003, citado en
\citeproc{ref-iturrasanhueza_estatus_2019}{Iturra Sanhueza, 2019, p.
15}). \textbf{Pendiente: Goldthorpe es una cita secundaria, a través de
Iturra. revisar gold nuevamente}

En relación a la educación se plantean dos hipótesis opuestas, la
primera es la hipótesis de la socialización que plantea que el colegio
transmite la idea de que el éxito depende del mérito individual, por lo
que quienes han avanzado más en el sistema educativo adherirían más a la
meritocracia (\citeproc{ref-bourdieu_reproduccion_1998}{Bourdieu \&
Passeron, 1998}; \citeproc{ref-lampert_meritocratic_2013}{Lampert,
2013}). Desde esta perspectiva, la escuela traduce el privilegio social
en talento, don o mérito
(\citeproc{ref-grau-grau_ilusion_2022}{Grau-Grau, 2022, p. 5}), y
quienes tienen mayor capital cultural atribuyen en mayor medida su
posición a su propio talento y esfuerzo
(\citeproc{ref-grau-grau_ilusion_2022}{Grau-Grau, 2022, p. 15}). La
hipótesis de la iluminación plantea, en cambio, que la educación aumenta
la sensibilidad frente a las desigualdades y las limitaciones sociales,
por lo que las personas más educadas serían más críticas respecto de
cómo se distribuyen las recompensas
(\citeproc{ref-bucca_merit_2016}{Bucca, 2016};
\citeproc{ref-duru-bellat_variations_2012}{Duru-Bellat \& Tenret,
2012}).

En América Latina, quienes tienen título universitario explican la
riqueza con más causas estructurales
(\citeproc{ref-bucca_merit_2016}{Bucca, 2016, p. 105}), y en Chile un
mayor nivel educacional se asocia con una menor percepción de
meritocracia (\citeproc{ref-castillo_merit_perc_pref_2019}{Castillo
et~al., 2019, p. 9};
\citeproc{ref-iturra-sanhueza_social_2026}{Iturra-Sanhueza et~al.,
2026}). Esto se ha explicado porque los más educados tienen más contacto
con grupos socioeconómicos diversos y porque sus experiencias de
movilidad no siempre se traducen en los retornos esperados
(\citeproc{ref-iturra-sanhueza_social_2026}{Iturra-Sanhueza et~al.,
2026, p. 9}). Sin embargo, ambas hipótesis pueden operar a la vez. En
Chile, las personas más educadas atribuyen más importancia a la
ambición, la educación y el trabajo duro para salir adelante, como
predice la socialización, pero perciben en menor medida que el esfuerzo
y la inteligencia se recompensan, como predice la iluminación
(\citeproc{ref-iturrasanhueza_estatus_2019}{Iturra Sanhueza, 2019, pp.
41-42, 45}).

\textbf{Iturra-Sanhueza et al.~2026 p.~9 se apoya en Londoño-Vélez,
Espinoza y Núñez y Canales, así que es una cita secundaria. revisar para
citar bien y leer a londoño}

La evidencia que distingue entre ambos criterios es escasa y no apunta
en una sola dirección. Con la medida de recompensa percibida, la
educación se asocia negativamente tanto con el esfuerzo como con el
talento, con magnitudes similares
(\citeproc{ref-iturra-sanhueza_social_2026}{Iturra-Sanhueza et~al.,
2026}, Tabla S3). En estudiantes chilenos, la educación universitaria de
los padres se asocia con una menor importancia del talento para obtener
buenas notas, pero no del esfuerzo
(\citeproc{ref-retamales_rol_2024}{Retamales, 2024}), y en Flandes los
más educados recurren menos al talento para explicar el éxito escolar
(\citeproc{ref-spruyt_talent_esfuerzo_social_background_2015}{Spruyt,
2015, p. 104}). En cambio, las personas con educación universitaria
prefieren que las habilidades se recompensen más y el esfuerzo menos
(\citeproc{ref-kwon_access_2026}{Kwon \& Erola, 2026, p. 1010}), aunque
se trata de preferencias y no de percepciones.

\textbf{{[}pendiente: Retamales, agregar la pagina. Spruyt p.~104. Kwon
y Erola p.~1010.{]}}

Esto sugiere que la iluminación y la socialización podrían operar de
forma distinta según el criterio. El esfuerzo es lo que la experiencia
educativa muestra como insuficiente, ya que los más educados conocen
mejor las barreras que impiden que se traduzca en recompensas. El
talento, en cambio, es lo que la escuela certifica al traducir las
ventajas de origen en capacidades individuales
(\citeproc{ref-accominotti_making_2026}{Accominotti \& Sauder, 2026}).
Por ello se espera que la educación se asocie negativamente con la
percepción de que se recompensa el esfuerzo y positivamente con la de
que se recompensa el talento.

\textbf{{[}pendiente: las dos oraciones sobre el esfuerzo y el talento
son una síntesis de las citas anteriores. Revisarlas con p JC{]}}

\begin{quote}
\(H_{2a}\): A mayor nivel educacional, menor es la percepción de que en
Chile se recompensa el esfuerzo.

\textbf{\textgreater{} \(H_{2b}\): A mayor nivel educacional, mayor es
la percepción de que en Chile se recompensa el talento.}

\(H_{2c}\): El nivel educacional se asocia de forma más positiva con la
percepción de que se recompensa el talento que con la de que se
recompensa el esfuerzo.
\end{quote}

\subsubsection{Movilidad social
subjetiva}\label{movilidad-social-subjetiva}

A diferencia de la movilidad social objetiva que refiere a los cambios
que experimentan las personas en la sociedad en relación con la
estructura socioeconómica, en cambio, la movilidad social subjetiva,
refiere a como las personas evalúan su posición actual en comparación
con la de la familia con la que crecieron. La movilidad subjetiva
permite comprender mejor qué actitudes y comportamientos se ven
influidos por la movilidad social, así como las razones por las cuales
esta los afecta (\citeproc{ref-kelley_subjective_2008}{Kelley \& Kelley,
2008}). Se considera la medida subjetiva y no la objetiva porque la
movilidad incide en las creencias en la medida en que las personas la
experimentan como real y tiene mayor relación con las experiencias de
las personas (\citeproc{ref-mijs_belief_2022}{Mijs et~al., 2022, p.
138}).

Para esta relación se han propuesto dos mecanismos. Por un lado, las
personas infieren cómo funciona la sociedad a partir de su propia
trayectoria. Por otro, quienes ascienden tienen motivos para presentar
su éxito como merecido (\citeproc{ref-mijs_belief_2022}{Mijs et~al.,
2022, p. 133}). Las personas tienden a generalizar su experiencia de
movilidad al conjunto de la sociedad
(\citeproc{ref-traini_et_al_mov_social_percepciones_meritocracia_2025}{Traini
et~al., 2025, p. 4}), y quienes logran ascender tienden a atribuir su
éxito a su propio trabajo y habilidad, y a usar esa experiencia para
explicar el éxito o el fracaso de los demás
(\citeproc{ref-bucca_merit_2016}{Bucca, 2016, p. 99}).

Para el esfuerzo, la evidencia internacional respalda esta relación. En
Países Bajos, quienes perciben haber ascendido atribuyen más importancia
al trabajo duro para salir adelante que quienes se mantuvieron en la
misma posición (\citeproc{ref-mijs_belief_2022}{Mijs et~al., 2022, p.
136}). Con datos de 26 países, entre ellos Chile, quienes perciben haber
ascendido consideran que el trabajo duro es más importante para tener
éxito
(\citeproc{ref-traini_et_al_mov_social_percepciones_meritocracia_2025}{Traini
et~al., 2025, p. 11}).

Sin embargo, estos estudios miden la importancia atribuida al trabajo
duro para salir adelante y no la percepción de que efectivamente se
recompensa, dos dimensiones que en Chile se comportan como factores
distintos (\citeproc{ref-iturrasanhueza_estatus_2019}{Iturra Sanhueza,
2019, p. 35}). \textbf{El problema es que cuando se mide la recompensa
percibida, la evidencia chilena no es clara.} A igual estatus subjetivo
propio, quienes ubican a su familia de origen en una posición más alta
perciben en mayor medida que se recompensan el esfuerzo y la
inteligencia (\citeproc{ref-iturrasanhueza_estatus_2019}{Iturra
Sanhueza, 2019, p. 46}). No obstante, ese resultado proviene de una
medida que combina ambos criterios y trata el origen como una escala
continua, por lo que no permite saber si la asociación es la misma para
el esfuerzo y para el talento, ni si difiere entre quienes ascienden y
quienes descienden. En América Latina, la movilidad percibida respecto
de los padres no se asocia con las explicaciones de la riqueza y la
pobreza (\citeproc{ref-bucca_merit_2016}{Bucca, 2016, pp. 105-106}), y
en Chile quienes han tenido una movilidad ascendente de largo alcance
explican su trayectoria desde un marco colectivo y reconocen los
factores externos que la hicieron posible, lo que menciona como
conciencia dual (\citeproc{ref-fercovic_disentanglin_2022}{Fercovic,
2022}).

\textbf{{[}pendiente Iturra p.~35 y Tabla 8 en p.~46. Confirmar que un
valor alto en ``recompensa percibida'' significa más recompensa (57),
que dice lo contrario. bucca pp.~105 y 106. fercovic.{]}}

La evidencia sobre el talento es más escasa y proviene de estudios
cualitativos sobre cómo las personas explican su propia trayectoria.
Quienes ascienden ponen el trabajo duro al centro de sus relatos. Entre
las élites de Reino Unido y Dinamarca, quienes ascendieron destacan más
el trabajo duro que el resto
(\citeproc{ref-friedman_meaning_2024}{Friedman et~al., 2024, p. 876}), y
estudiantes chinos de clase trabajadora que llegaron a universidades de
élite sostienen que el trabajo duro siempre da frutos
(\citeproc{ref-jin_meritocracy_2020}{Jin \& Ball, 2020, p. 10}). Pero
también recurren al talento. En Lima, profesionales que ascendieron
perciben que en los campos de alta calificación el talento y la
inteligencia pesan más que las ventajas de un origen privilegiado, y
\textbf{usan el talento para distinguirse de sus colegas privilegiados
(\citeproc{ref-renteria_telling_2026}{Rentería, 2026, p. 1353}).} Si las
personas proyectan su experiencia hacia la sociedad, esto sugiere que la
movilidad ascendente se asocia con ambas percepciones, pero de forma más
positiva con la del esfuerzo, ya que es el criterio que depende en mayor
medida de la voluntad
(\citeproc{ref-franetovic_concepciones_2017a}{Franetovic, 2017, p. 5}) y
que permite presentar el ascenso como resultado del trabajo propio. Dado
que la evidencia sobre la movilidad descendente es menos consistente
(\citeproc{ref-mijs_belief_2022}{Mijs et~al., 2022, p. 136};
\citeproc{ref-traini_et_al_mov_social_percepciones_meritocracia_2025}{Traini
et~al., 2025, pp. 11-12}), esta se incluye en los análisis sin una
hipótesis específica.

o \emph{acá hay dos argumentos contrarios que orientan las hipotesis,
tengo una duda acá, porque no sé si están correctos los argumentos}

Para el talento, la evidencia es más escasa y apunta en otra dirección.
En Lima, profesionales que ascendieron perciben que en los campos de
alta calificación el talento y la inteligencia pesan más que las
ventajas de un origen privilegiado, y usan el talento para distinguirse
de sus colegas privilegiados
(\citeproc{ref-renteria_telling_2026}{Rentería, 2026, p. 1353}). Además,
frente a sus éxitos, las personas tienden a exagerar su propia capacidad
y talento para reforzar una imagen positiva de sí mismas
(\citeproc{ref-traini_et_al_mov_social_percepciones_meritocracia_2025}{Traini
et~al., 2025, p. 3}). Esto sugiere que, mientras su trayectoria les
muestra los límites de la recompensa al esfuerzo, el talento les permite
explicar por qué ellos sí lograron ascender. Como esta evidencia
proviene de un solo estudio cualitativo, la hipótesis para el talento es
más tentativa. Dado que la evidencia sobre la movilidad descendente es
menos consistente (\citeproc{ref-mijs_belief_2022}{Mijs et~al., 2022, p.
136};
\citeproc{ref-traini_et_al_mov_social_percepciones_meritocracia_2025}{Traini
et~al., 2025, pp. 11-12}), esta se incluye en los análisis sin una
hipótesis específica.

\begin{quote}
\(H_{3a}\): A igual estatus social subjetivo, quienes perciben haber
ascendido respecto de su familia de origen perciben en mayor medida que
en Chile se recompensa el esfuerzo que quienes perciben mantenerse en la
misma posición.

\(H_{3b}\): \textbf{A igual estatus social subjetivo, quienes perciben
haber ascendido respecto de su familia de origen perciben en mayor
medida que en Chile se recompensa el talento que quienes perciben
mantenerse en la misma posición.}

\(H_{3c}\): La movilidad ascendente se asocia de forma más positiva con
la percepción de que se recompensa el esfuerzo que con la de que se
recompensa el talento.
\end{quote}

\subsubsection{Nivel Socioeconómico}\label{nivel-socioeconuxf3mico}

El nivel socioeconómico refiere a la posición objetiva que ocupan las
personas en la estructura social según sus recursos
(\citeproc{ref-li_community_2021}{Li \& Hu, 2021, p. 478}). En este
estudio se aborda a partir del ingreso del hogar, que junto con la
educación conforma el estatus objetivo
(\citeproc{ref-iturra-sanhueza_social_2026}{Iturra-Sanhueza et~al.,
2026, p. 5}), con el fin de identificar a quienes ocupan las posiciones
más altas. El estatus social subjetivo no permite hacerlo, ya que en
Chile la autoubicación tiende hacia el centro, sobre todo entre quienes
tienen un mayor estatus objetivo
(\citeproc{ref-castillo_todos_clase_media_2013}{Castillo et~al., 2013}).
Así, quienes están en la parte alta de la distribución suelen ubicarse a
sí mismos en posiciones medias, por lo que su forma de percibir la
meritocracia podría quedar oculta si solo se considera el estatus
subjetivo.

\textbf{{[}pendiente: li y Hu, p.~478 (leer nuevamente). turra-Sanhueza
et al.~2026, p.~5. Castillo et al.~2013:revisar la página.{]}}

La evidencia chilena indica que quienes ocupan las posiciones más altas
explican el acceso a ellas por el talento más que por el esfuerzo. La
élite económica ancla el mérito en el talento, a diferencia de las
clases medias y bajas, que lo equiparan con el esfuerzo
(\citeproc{ref-atria_et_al_elites_meritocracia_2020}{Atria et~al., 2020,
p. 3}). Más de la mitad de sus integrantes asocia explícitamente el
mérito con el talento o con conceptos similares, como habilidades o
dones (\citeproc{ref-atria_et_al_elites_meritocracia_2020}{Atria et~al.,
2020, p. 10}), de modo que se privilegia el talento por sobre el
esfuerzo (\citeproc{ref-atria_et_al_elites_meritocracia_2020}{Atria
et~al., 2020, p. 16}).

\textbf{{[}Pendiente: atria pp.~3, 10 y 16¿ paginación de la revista es
1219 a 1241. Quaresma: confirmar la llave de las dos investigaciones q
tiene {]}}

Fuera de Chile sucede algo parecido, en Reino Unido las élites destacan
su talento natural, una idea que cultivan las escuelas privadas de
élite, mientras que en Dinamarca destacan el trabajo duro
(\citeproc{ref-friedman_meaning_2024}{Friedman et~al., 2024, p. 861}).
Además, los estudiantes de clase media y de élite tienden a verse como
naturalmente inteligentes, mientras que los de clase trabajadora deben
demostrar su esfuerzo para no ser vistos como perezosos
(\citeproc{ref-jackson_smart_2015}{Jackson \& Nyström, 2015, p. 403}).

\textbf{Friedman, p.~861 p.~871. Jackson y Nyström revisar paginas}

Pero, en Chile el ingreso no muestra una asociación lineal con la
percepción de que se recompensa el esfuerzo ni con la de que se
recompensa el talento
(\citeproc{ref-iturra-sanhueza_social_2026}{Iturra-Sanhueza et~al.,
2026}, Tabla S3), pero el 10\% de mayores ingresos percibe más
meritocracia que el resto
(\citeproc{ref-iturra-sanhueza_social_2026}{Iturra-Sanhueza et~al.,
2026}, Tabla S1). Como esa medida combina ambos criterios, no permite
saber si la diferencia proviene del esfuerzo, del talento o de ambos.
Aun así, a partir de la evidencia sobre las élites, se espera que los
hogares de mayores ingresos perciban en mayor medida que se recompensa
el talento y en menor medida que se recompensa el esfuerzo. Los hogares
de altos ingresos no equivalen a la élite económica, que se define por
ocupar posiciones de poder en grandes empresas
(\citeproc{ref-atria_et_al_elites_meritocracia_2020}{Atria et~al.,
2020}). Sin embargo, la asociación entre mérito y talento no se limita a
la élite, ya que, también aparece entre directivos de universidades de
élite (\citeproc{ref-quaresma_merit_u_elites_2024}{Quaresma, Madrid,
et~al., 2024}) y entre estudiantes de clase media y de élite, que
tienden a verse como naturalmente inteligentes
(\citeproc{ref-jackson_smart_2015}{Jackson \& Nyström, 2015, p. 403}).
Por ello se espera que los hogares de altos ingresos compartan, al menos
en parte, esta forma de entender el mérito.

\textbf{{[}Pendiente: Tablas S1 y S3 del material suplementario de
iturra-Ssanhueza et al.~2026.{]}}

\begin{quote}
\(H_{4a}\): Quienes pertenecen a hogares de altos ingresos perciben en
menor medida que en Chile se recompensa el esfuerzo que quienes
pertenecen a hogares de ingresos medios y bajos.

\(H_{4b}\): Quienes pertenecen a hogares de altos ingresos perciben en
mayor medida que en Chile se recompensa el talento que quienes
pertenecen a hogares de ingresos medios y bajos.

\(H_{4c}\): La diferencia entre los hogares de altos ingresos y el resto
es más positiva para la percepción de que se recompensa el talento que
para la de que se recompensa el esfuerzo.
\end{quote}

\subsubsection{Variables de control}\label{variables-de-control}

El sexo porque las ``También incluimos el sexo, dada la evidencia de que
las mujeres tienden a favorecer la protección estatal más que los
hombres, en parte debido a trayectorias diferenciadas del mercado
laboral'' Plaza et~al. (\citeproc{ref-plaza_networked_2026}{2026})

\section{Resumen del capítulo}\label{resumen-del-capuxedtulo}

\begin{quote}
\emph{Este capítulo ofreció una revisión breve de los conceptos
fundamentales y estableció proposiciones para guiar el análisis empírico
posterior} \textbf{cambiar, es de la plantilla facso}
\end{quote}

\bookmarksetup{startatroot}

\chapter{Metodología}\label{metodologuxeda}

\section{Datos}\label{datos}

En el presente estudio, se utilizarán los datos de la encuesta
``Economía Moral de la Meritocracia y Preferencias Redistributivas'' una
encuesta parte del proyecto Fonfecyt 1160921 llevada a cabo en abril
2016 - marzo 2020. Los análisis se basan en datos de la primera ola,
recopilados en \emph{año}. Esta encuesta tiene una muestra de N = 2.457.
Y, contempla un cuestionario de 37 preguntas aplicado mediante
Qualtrics. Los encuestados responden preguntas sobre temas como
tolerancia a la desigualdad, importancia para surgir en la vida,
igualitarismo, percepciones y preferencias en meritocracia, entre otros.

La encuesta tiene dos puntos fuertes en relación con la investigación de
meritocracia. Incluye percepciones y preferencias meritocráticas y junto
a este, factores no meritocráticos, lo que permite dimensionar el
concepto más allá de las creencias y también separa el concepto en sus
dos dimensiones, es decir, talento y esfuerzo.

\section{Variables dependientes: Percepciones de esfuerzo y
talento}\label{variables-dependientes-percepciones-de-esfuerzo-y-talento}

Las variables dependientes del estudio son las percepciones de que se
premia el esfuerzo y el talento en la sociedad por separado, ambas
percepciones se miden con un ítem tipo Likert de cinco categorías. La
dimensión del esfuerzo, se basa en la afirmación: 'Quienes más se
esfuerzan logran obtener mayores recompensas que quienes se esfuerzan
menos'. Medida en una escala Likert de 5 puntos, desde 1 (totalmente en
desacuerdo) hasta 5 (totalmente de acuerdo). Y la dimensión del talento,
se basa en la afirmación: 'Quienes poseen más talento logran obtener
mayores recompensas que quienes poseen menos talento'. Medida en una
escala Likert de 5 puntos, desde 1 (totalmente en desacuerdo) hasta 5
(totalmente de acuerdo). Los valores altos en las variables indican
percepciones afirmativas sobre que se recompensa el talento y esfuerzo
en la sociedad.

\section{Variables independientes}\label{variables-independientes}

Para las variables independientes del estudio se utilizará el estatus
social subjetivo que se mide como una escala de autoubicación de 11
puntos donde se pregunta: `En Chile, hay grupos que tienden a ubicarse
en los niveles más altos y grupos que tienden a ubicarse en los niveles
más bajos de la sociedad. Usando la escala presentada, donde 0 es el
nivel más bajo y 10 el nivel más alto, ¿dónde se ubicaría usted en la
sociedad chilena?'. En la cual valores más altos significan mayor
estatus social percibido.

El nivel educacional se mide como una escala de autoubicación de 10
categorías del nivel que se cursa actualmente o el último aprobado,
donde se pregunta: `¿Cuál es su mayor nivel educacional alcanzado?' La
escala va desde (1) No estudió, (2) Educación básica incompleta, (3)
Educación básica completa, (4) Educación media incompleta, (5)Educación
media completa, (6) Educación superior no universitaria incompleta, (7)
Educación superior no universitaria completa, (8) Educación
universitaria incompleta, (9) Educación universitaria completa, y (10)
Estudios de postgrado, master, doctorado. Donde valores más altos
significan mayor nivel educacional alcanzado. \textbf{ver si codifico
según oit, mctp o CINE antigua de 1997}

La movilidad social subjetiva se construye con la variable de estatus
subjetivo propio y el de la familia de origen, el ítem menciona: `\,``Y
pensando en la familia en la que usted creció, ¿dónde se ubicarían ellos
en esta escala?'', de 0 a 10'. Siguiendo a Mijs et~al.
(\citeproc{ref-mijs_belief_2022}{2022, p. 134}) y Kwon \& Erola
(\citeproc{ref-kwon_access_2026}{2026, p. 1008}), ambas escalas se
agruparán en bajo (0 a 3), medio (4 a 6) y alto (7 a 10) y se cruzan.
Quienes quedan en un nivel superior al de su familia se clasifican como
ascendentes, quienes quedan en uno inferior como descendentes y quienes
se mantienen como inmóviles.

El nivel socioeconómico se mide a través del ingreso del hogar en una
escala de 14 puntos de tramos, donde se pregunta: `De los siguientes
tramos de ingresos mensuales que se presentan, ¿podría Ud. indicarme en
cuál de ellos se encuentra su hogar, considerando todos los ingresos
líquidos por sueldos y salarios de todas las personas que trabajan
remuneradamente, jubilaciones, pensiones, aportes de parientes o amigos,
arriendos y otros?'. La escala va desde 1, la cual representa ingresos
menores a \$35.000 mensuales líquidos, al 14, que representa más de
\$3.000.000 mensuales líquidos. La variable será recodificada siguiendo
la estrategia de estimar diferencias entre grupos de ingreso en lugar de
un coeficiente único (\citeproc{ref-mac_baroz_merit_des_2024}{Mac-Clure
et~al., 2024, p. 13}), porque no se puede construir quintiles de ingreso
per cápita al no preguntar por el número de integrantes del hogar. Los
tramos se agrupan en cuatro categorías: ingresos bajos (hasta \$448.000,
referencia), medio-bajos (\$448.001 a \$1.000.000), medio-altos
(\$1.000.001 a \$2.000.000) y altos (más de \$2.000.000). Esta última se
aproxima al 10\% de mayores ingresos que distinguen Iturra-Sanhueza
et~al. (\citeproc{ref-iturra-sanhueza_social_2026}{2026}, Tabla S1).

\section{Variables de control}\label{variables-de-control-1}

La variable de sexo es dicotómica donde (1) es hombre y (2) es mujer. La
edad se utilizará como una variable continua.

\section{Métodos}\label{muxe9todos}

\textbf{Breve descripción de los procedimientos aplicados para procesar
y analizar la información.}

Se estima una regresión lineal multivariante en la que ambas
percepciones son variables dependientes simultáneas, con las mismas
variables independientes. Esto permite estimar si los factores asociados
de ambas variables dependientes afectan de igual manera. Test de Wald.
Los análisis se realizan en R.

\bookmarksetup{startatroot}

\chapter*{Prefacio}\label{prefacio}
\addcontentsline{toc}{chapter}{Prefacio}

\markboth{Prefacio}{Prefacio}

Agradecimientos opcionales, dedicatorias, notas sobre financiamiento,
etc.

\mainmatter

\bookmarksetup{startatroot}

\chapter{Introducción}\label{introducciuxf3n}

\section{Problema de
investigación}\label{problema-de-investigaciuxf3n-1}

La meritocracia es el principio de justicia más extendido en las
sociedades contemporáneas, entregando una base moral para asignar
recursos y recompensas escasas de manera proporcional al desempeño
individual, dejando de lado el origen social
(\citeproc{ref-mijs_unfulfillable_2016}{Mijs, 2016};
\citeproc{ref-mulligan_meritocracy_2023a}{Mulligan, 2023}). Además,
cumple la función de legitimar la desigualdad, ya que al reducir el
éxito a factores individuales deja de lado los factores estructurales y
presenta las brechas sociales como merecidas
(\citeproc{ref-accominotti_making_2026}{Accominotti \& Sauder, 2026};
\citeproc{ref-castillo_merit_perc_pref_2019}{Castillo et~al., 2019};
\citeproc{ref-lierse_des_redi_2022}{Lierse et~al., 2022, p. 373}). Desde
la sociología, esta legitimación se ha estudiado a través de las
percepciones meritocráticas, es decir, lo que las personas creen sobre
si la sociedad efectivamente recompensa el esfuerzo y el talento
(\citeproc{ref-castillo_merit_perc_pref_2019}{Castillo et~al., 2019}).
Estas percepciones varían según la posición que las personas ocupan y la
trayectoria que han recorrido, en particular según su estatus social
subjetivo, su nivel educacional y su experiencia de movilidad social
(\citeproc{ref-iturra-sanhueza_social_2026}{Iturra-Sanhueza et~al.,
2026}; \citeproc{ref-mijs_belief_2022}{Mijs et~al., 2022}). Sin embargo,
los factores asociados a percibir que se recompensa el esfuerzo podrían
no ser los mismos que los asociados a percibir que se recompensa el
talento.

El esfuerzo y el talento de la fórmula de Young
(\citeproc{ref-young_1964}{1964}) contienen dos criterios que remiten a
principios de merecimiento distintos
(\citeproc{ref-clavero_2023}{Clavero García, 2023};
\citeproc{ref-sardoc_rise_2024}{Sardoč, 2024}). El esfuerzo depende en
mayor medida de la voluntad y la motivación de las personas
(\citeproc{ref-franetovic_concepciones_2017a}{Franetovic, 2017, p. 5}),
mientras que el talento, suele entenderse como un conjunto de
capacidades innatas distribuidas de manera azarosa
(\citeproc{ref-harden_2021}{Harden, 2021};
\citeproc{ref-wright_preguntas_2010}{Wright, 2010, pp. 50-51}), escapa
en mayor medida del control de las personas.

Las personas recurren al mérito de manera discursiva para explicar,
justificar y dar sentido a sus posiciones en la estructura social. De
esta manera, recurrir al mérito permite legitimar la propia posición,
aun cuando se sostiene en ventajas heredadas
(\citeproc{ref-khan_saying_2013}{Khan \& Jerolmack, 2013, p. 12}). En
Chile, las élites económicas explican el acceso a las posiciones
superiores más por el talento que por el esfuerzo, mientras que en las
clases medias y bajas, y en las universidades de masas, predomina el
esfuerzo (\citeproc{ref-atria_et_al_elites_meritocracia_2020}{Atria
et~al., 2020}; \citeproc{ref-quaresma_merit_u_elites_2024}{Quaresma,
Madrid, et~al., 2024}). Quienes han experimentado movilidad ascendente
tienden a destacar su esfuerzo (\citeproc{ref-jin_meritocracy_2020}{Jin
\& Ball, 2020}; \citeproc{ref-mijs_belief_2022}{Mijs et~al., 2022}),
aunque también pueden reivindicar su talento para distinguirse de sus
pares privilegiados (\citeproc{ref-renteria_telling_2026}{Rentería,
2026}), y el criterio que privilegian las élites varía según el país y
las instituciones en las que se forman
(\citeproc{ref-friedman_meaning_2024}{Friedman et~al., 2024};
\citeproc{ref-gogescu_meritocracy_2025}{Gogescu, 2025};
\citeproc{ref-warikoo_legitimating_2014}{Warikoo \& Fuhr, 2014}).

Hasta ahora, se han analizado las percepciones meritocráticas bajo la
fórmula de Young (\citeproc{ref-young_1964}{1964}), pero rara vez se
analiza el talento por separado\footnote{La encuesta ISSP (International
  Social Survey Programme) incluye en su módulo de desigualdad social la
  importancia de las habilidades naturales para salir adelante: ``Get
  ahead: How important is having natural ability?''
  (\citeproc{ref-issp_data}{ISSP Research Group, 2002}). ELSOC (Estudio
  Longitudinal Social de Chile) incluye el ítem ``En Chile, las personas
  son recompensadas por su inteligencia y habilidad''
  (\citeproc{ref-elsoc_data}{Reproducible Research, Centre for Social
  Conflict and Cohesion Studies COES, 2024}). La encuesta CFPS (China
  Family Panel Studies) pregunta ``The most important factor affecting
  one's future success is his/her talent'' y, en olas enfocadas en el
  rendimiento educativo, ``In your opinion, talent accounts for
  \_\_\_\_\_\_\% in achieving a good academic record''
  (\citeproc{ref-cfps_data}{Institute of Social Science Survey, Peking
  University, 2015}).}. Algunos estudios las miden solo a través del
esfuerzo (\citeproc{ref-ehmes_paradox_2026}{Ehmes \& Gangl, 2026};
\citeproc{ref-mijs_paradox_2021}{Mijs, 2021};
\citeproc{ref-wiesner_effort_2026}{Wiesner \& Sachweh, 2026}) y otros
suman esfuerzo y talento en un mismo indicador
(\citeproc{ref-castillo_et_al_multidimensional_merit_beliefs_2023}{Castillo
et~al., 2023};
\citeproc{ref-iturra-sanhueza_social_2026}{Iturra-Sanhueza et~al.,
2026}). Así, los factores asociados a la percepción meritocrática se han
estimado sobre medidas que no distinguen entre ambos criterios, y no es
posible saber si los resultados valen para los dos o solo para uno de
ellos.

Hay indicios de que ambos criterios se comportan de manera distinta. En
Chile, el esfuerzo se percibe como menos recompensado que la
inteligencia y las habilidades
(\citeproc{ref-castillo_merit_perc_pref_2019}{Castillo et~al., 2019, p.
6}; \citeproc{ref-iturra-sanhueza_social_2026}{Iturra-Sanhueza et~al.,
2026, p. 6}), y cuando ambas percepciones se estiman por separado, su
asociación con el estatus subjetivo y la educación difiere en magnitud,
aunque esas diferencias no han sido puestas a prueba formalmente
(\citeproc{ref-iturra-sanhueza_social_2026}{Iturra-Sanhueza et~al.,
2026}, Tabla S3). En este escenario, el presente estudio busca
establecer en qué medida los factores asociados a la percepción de que
en Chile se recompensa el esfuerzo difieren de aquellos asociados a la
percepción de que se recompensa el talento.

Para ello, se consideran cuatro factores asociados a la posición y a la
trayectoria social de las personas. \emph{Primero}, el nivel
socioeconómico, ya que quienes ocupan posiciones objetivamente más
aventajadas tienden a explicar su lugar por el talento
(\citeproc{ref-atria_et_al_elites_meritocracia_2020}{Atria et~al.,
2020}) y el 10\% de mayores ingresos percibe más meritocracia que el
resto (\citeproc{ref-iturra-sanhueza_social_2026}{Iturra-Sanhueza
et~al., 2026}). El \emph{segundo} es el estatus social subjetivo, es
decir, cómo las personas se autoidentifican en la jerarquía social, ya
que quienes se perciben en posiciones más altas tienden a percibir que
la sociedad premia el mérito, lo que legitima su propia posición
(\citeproc{ref-castillo_merit_perc_pref_2019}{Castillo et~al., 2019};
\citeproc{ref-iturra-sanhueza_social_2026}{Iturra-Sanhueza et~al.,
2026}). El \emph{tercero} es el nivel educacional, ya que la educación
se asocia con una visión más crítica de cómo se distribuyen las
recompensas, pero al mismo tiempo certifica los logros académicos a
través del mérito individual
(\citeproc{ref-castillo_merit_perc_pref_2019}{Castillo et~al., 2019};
\citeproc{ref-iturra-sanhueza_social_2026}{Iturra-Sanhueza et~al.,
2026}). El cuarto es, la movilidad social subjetiva, ya que quienes
perciben haber ascendido tienden a interpretar la sociedad a partir de
su propia trayectoria y a percibirla como más meritocrática
(\citeproc{ref-mijs_belief_2022}{Mijs et~al., 2022};
\citeproc{ref-traini_et_al_mov_social_percepciones_meritocracia_2025}{Traini
et~al., 2025}).

\textbf{ultima frase para resultados esperados de las hipotesis}

\section{Objeto de estudio}\label{objeto-de-estudio-1}

La investigación sobre meritocracia distingue entre las preferencias,
que son creencias normativas sobre cómo debería funcionar la sociedad, y
las percepciones, que son descripciones de cómo funciona y que pueden
diferir de lo que efectivamente sucede en ella
(\citeproc{ref-castillo_percepcion_2012}{Castillo et~al., 2012};
\citeproc{ref-castillo_merit_perc_pref_2019}{Castillo et~al., 2019}).
Percibir que una sociedad premia el esfuerzo o el talento es, por lo
tanto, distinto de preferir que lo haga. Según el teorema de Thomas, si
las personas definen una situación como real, esta es real en sus
consecuencias (\citeproc{ref-merton_selffulfilling_1948}{Merton, 1948}).
Así, percibir que la sociedad es meritocrática se asocia con considerar
más aceptables las desigualdades (\citeproc{ref-dubet_nuevo_2023}{Dubet,
2023, pp. 32-33}) y con un menor apoyo a la redistribución
(\citeproc{ref-sonhing_merit_2011}{Son Hing et~al., 2011};
\citeproc{ref-tejero-peregrina_perceived_2025}{Tejero-Peregrina et~al.,
2025}). Además, estas consecuencias difieren según el criterio. En
Chile, percibir que se recompensa el esfuerzo se asocia con legitimar
que el acceso a la salud, la educación y las pensiones dependa del
ingreso, mientras que percibir que se recompensa el talento no se
asocia, o se asocia negativamente, con esa legitimación
(\citeproc{ref-castillo_perceptions_2025}{Castillo et~al., 2025},
\citeproc{ref-castillo_justification_desigualdades_pensiones_clases_2026}{2026}).

Así, esta investigación tiene como objeto de estudio las percepciones
sobre la meritocracia, diferenciando entre la percepción de que se
recompensa el esfuerzo y la percepción de que se recompensa el talento.

\section{Pregunta de
investigación}\label{pregunta-de-investigaciuxf3n-1}

¿En qué medida el nivel socioeconómico, el estatus social subjetivo, el
nivel educacional y la experiencia de movilidad social se asocian con la
percepción de que en Chile se recompensa el esfuerzo y con la percepción
de que se recompensa el talento, y en qué medida difieren estas
asociaciones entre ambas percepciones?

\bookmarksetup{startatroot}

\chapter{Antecedentes conceptuales y
empíricos}\label{antecedentes-conceptuales-y-empuxedricos}

\section{Marco conceptual}\label{marco-conceptual-1}

\subsection{Meritocracia: concepto y
debate}\label{meritocracia-concepto-y-debate-1}

En la sociología el mérito es un tipo de merecimiento propio de la
justicia distributiva, que depende de las normas, valores y objetivos de
prácticas sociales particulares (\citeproc{ref-pavel_merit_2024}{Pavel,
2024}). El concepto de meritocracia se populariza con Young
(\citeproc{ref-young_1964}{1964}) como una sátira a una sociedad
gobernada por élites, sustentada en la fórmula mérito = inteligencia +
esfuerzo (M = I + E). La idea supone la igualdad de oportunidades, en la
cual los individuos, a través del esfuerzo y el talento, pueden lograr
movilidad social y éxito en ciertas esferas sociales, dejando de lado
factores no meritocráticos como el género, la clase social o el origen
familiar (\citeproc{ref-wooldridge_aristocracy_2021}{Wooldridge, 2021}).
El concepto ofrece así, una base moral para la asignación de bienes
escasos, ya que recompensa lo que las personas hacen, en proporción a su
desempeño (\citeproc{ref-mijs_unfulfillable_2016}{Mijs, 2016}). Según
Mijs (\citeproc{ref-mijs_unfulfillable_2016}{2016}), la meritocracia se
compone de tres elementos fundamentales: las carreras abiertas al
talento, las oportunidades educativas adaptadas a la capacidad natural y
el logro, en lugar de la adscripción, como base de la desigualdad
social.

El contenido del mérito es difuso y depende de lo que cada sociedad
valora. Para Sen (\citeproc{ref-sen_merit_justice_2000}{2000}), el
término está subdefinido y no hay claridad sobre el significado de
meritocracia. Premiar el mérito significa premiar lo que una sociedad
considera éxito, según su concepción de una buena sociedad, lo que
permitiría un procedimiento legítimo para asignar recursos y
recompensas. Por eso, para determinar lo que es meritorio y lo que no,
se necesita un grado de acuerdo sobre principios de justicia. En una
línea similar, Clavero García (\citeproc{ref-clavero_2023}{2023, p.
1180}) define el mérito como ``the value of a free, effortful,
good-oriented action'' {[}el valor de una acción libre, esforzada y
orientada al bien{]}. Como lo que se considera un bien varía entre
sociedades, el mérito sería contextual. En su formulación clásica, el
mérito combina talento y esfuerzo, entendido este último como la
dedicación y el trabajo sostenido hacia un fin
(\citeproc{ref-clavero_2023}{Clavero García, 2023};
\citeproc{ref-franetovic_concepciones_2017a}{Franetovic, 2017}). Bajo la
definición de Clavero, en cambio, el talento sólo sería meritorio en la
medida en que se ejerce mediante una acción libre y esforzada, y no por
sí mismo, lo que tensiona la fórmula de Young
(\citeproc{ref-young_1964}{1964}), que lo incluye como componente
directo del mérito.

Una primera distinción separa la meritocracia como ideal, es decir, cómo
deberían distribuirse las recompensas, de la meritocracia como
descripción de cómo funciona efectivamente la sociedad. Como ideal, la
meritocracia se defiende porque permite juzgar a las personas por sus
méritos y no por las ventajas heredadas, el nepotismo o el clientelismo,
algo que han exigido los sectores históricamente desfavorecidos, por lo
que sus problemas se resolverían profundizándola
(\citeproc{ref-wooldridge_aristocracy_2021}{Wooldridge, 2021}). La
crítica, en cambio, sostiene que el problema no es sólo que la
meritocracia esté incompleta, sino el principio mismo, ya que hacer
creer que las posiciones son resultado del esfuerzo y el talento
individual genera soberbia en quienes triunfan, humillación en quienes
quedan rezagados y convierte los títulos en una medida del valor de las
personas (\citeproc{ref-sandel_tirania_2020}{Sandel, 2020}).

El debate normativo no permite responder si la sociedad funciona de
manera meritocrática, por eso la pregunta iría orientada hacia si las
sociedades funcionan efectivamente de manera meritoria. Por un lado, la
promesa de la meritocracia basada en la educación no ha permitido un
aumento constante en la movilidad social de las clases trabajadoras,
según la cual las posiciones de clase dependerían cada vez más de las
credenciales educativas y menos del origen social, mostrando sólo
cambios leves y discontinuos
(\citeproc{ref-goldthorpe_myth_2003}{Goldthorpe, 2003, pp. 234-239}). De
esta manera, las oportunidades para tener mérito dependerían de factores
no meritocráticos, como los recursos familiares
(\citeproc{ref-mijs_unfulfillable_2016}{Mijs, 2016}).

Más allá de si la meritocracia se cumple, la sociología crítica
distingue entre el mérito como criterio de asignación y las ventajas de
origen, y plantea que el discurso meritocrático opera como un mecanismo
para legitimar el privilegio. El sistema educativo ejerce una violencia
simbólica al presentar como méritos individuales diferencias que
dependen del origen social, con lo que las desigualdades aparecen como
diferencias naturales entre las personas
(\citeproc{ref-bourdieu_reproduccion_1998}{Bourdieu \& Passeron, 1998}).
Este uso del discurso se ha descrito como decir meritocracia y hacer
privilegio, ya que estudiantes de élite explican su posición por el
trabajo duro, mientras en la práctica desarrollan una soltura que
protege las ventajas que heredaron (\citeproc{ref-khan_saying_2013}{Khan
\& Jerolmack, 2013}). Así, el trabajo duro opera como una cobertura
retórica de esas ventajas (\citeproc{ref-khan_saying_2013}{Khan \&
Jerolmack, 2013}). A su vez, la meritocracia no solo ordena a las
personas en posiciones desiguales, sino que construye órdenes de mérito,
es decir, jerarquías de estatus basadas en evaluaciones del mérito, el
logro o el desempeño, que establecen relaciones de superioridad,
igualdad o inferioridad entre las personas
(\citeproc{ref-accominotti_making_2026}{Accominotti \& Sauder, 2026}).
Estos órdenes moralizan la desigualdad al presentar las diferencias de
ventajas como diferencias de mérito individual. Además, pueden erosionar
la cohesión social y la solidaridad distributiva al reintroducir
jerarquías morales sobre el valor de las personas
(\citeproc{ref-accominotti_making_2026}{Accominotti \& Sauder, 2026};
\citeproc{ref-sandel_tirania_2020}{Sandel, 2020}), y existen como
esquemas culturales más o menos compartidos, de modo que lo que las
personas perciben como merecido forma parte de ellos
(\citeproc{ref-accominotti_making_2026}{Accominotti \& Sauder, 2026}).

En el estudio de las creencias meritocráticas se distinguen dos
enfoques, por un lado, las preferencias, que son creencias normativas
sobre cómo debería funcionar la sociedad, y por otro, las percepciones,
que son creencias descriptivas sobre cómo funciona efectivamente
(\citeproc{ref-castillo_merit_perc_pref_2019}{Castillo et~al., 2019}).
Como menciona Castillo et~al.
(\citeproc{ref-castillo_merit_perc_pref_2019}{2019}), percibir que la
meritocracia funciona no implica considerarla el mejor sistema ni
preferirla. Además, las percepciones cumplen una función legitimadora,
puesto que quienes perciben que la meritocracia funciona también
perciben menos desigualdad económica
(\citeproc{ref-castillo_merit_perc_pref_2019}{Castillo et~al., 2019, p.
11}). Estas percepciones no son descripciones neutras de la sociedad,
sino inferencias que las personas construyen a partir de su propia
trayectoria y de su entorno social, que permiten entender cómo es la
meritocracia en la sociedad (\citeproc{ref-dawtry_why_2015}{Dawtry
et~al., 2015}; \citeproc{ref-mijs_inequality_2018}{J. Mijs, 2018}).
Además, las personas reconocen factores meritocráticos y no
meritocráticos, como el origen familiar o los contactos, lo que se ha
denominado conciencia dual
(\citeproc{ref-reynolds_perceptions_2014}{Reynolds \& Xian, 2014};
\citeproc{ref-wiesner_effort_2026}{Wiesner \& Sachweh, 2026}). Por ello,
las percepciones meritocráticas y no meritocráticas se miden como
dimensiones distintas
(\citeproc{ref-castillo_et_al_multidimensional_merit_beliefs_2023}{Castillo
et~al., 2023}), y la forma en que se combinan varía entre países
(\citeproc{ref-zhu_meritocratic_2025}{Zhu, 2025}).

Si el discurso meritocrático permite legitimar la posición propia, la
percepción de que la sociedad recompensa el mérito no es independiente
del lugar que las personas ocupan en la estructura social, y como el
esfuerzo y el talento son criterios distintos, ese lugar podría
asociarse de manera diferente con cada uno.

\subsubsection{Talento y esfuerzo}\label{talento-y-esfuerzo}

El talento y el esfuerzo son los dos componentes del mérito en la
fórmula de Young (\citeproc{ref-young_1964}{1964}), y son también las
explicaciones más comunes del éxito y del fracaso
(\citeproc{ref-weiner_attributional_1985}{Weiner, 1985, p. 549}). Sin
embargo, no se sostienen en la misma base de merecimiento, ya que el
esfuerzo se percibe como algo que depende de las personas, mientras que
el talento ocupa una posición ambigua entre el logro y la herencia. Al
explicar el éxito escolar, por ejemplo, las personas en Flandes no
separan los factores meritocráticos de los no meritocráticos, sino lo
que puede manejar el individuo, el esfuerzo, de lo que viene dado, el
talento y el origen social
(\citeproc{ref-spruyt_talent_esfuerzo_social_background_2015}{Spruyt,
2015, p. 104}). Por ello, percibir que la sociedad recompensa el
esfuerzo implica creer que las recompensas dependen de lo que las
personas hacen, mientras que percibir que recompensa el talento implica
creer que dependen de lo que las personas tienen.

Además, las personas pueden percibir el esfuerzo y el talento de
distintas maneras, ya sea según su importancia para salir adelante, como
explicación de la pobreza y la riqueza, o como explicación de la propia
posición. Esta investigación se centra en la percepción de que la
sociedad recompensa a quienes poseen más talento y a quienes más se
esfuerzan (\citeproc{ref-castillo_merit_perc_pref_2019}{Castillo et~al.,
2019}; \citeproc{ref-castillo_perceptions_2025}{Castillo et~al., 2025}).
A continuación se analizan ambos criterios por separado.

\paragraph{Talento}\label{talento-1}

En esta investigación, el talento se entiende como una noción de sentido
común sobre las capacidades que las personas tienen, sin pronunciarse
sobre si estas son innatas o adquiridas. Por ello, interesa qué
reconocen las personas como talento cuando evalúan si la sociedad lo
recompensa. De esta manera, se presentan tres distinciones, si este es
innato o adquirido, si es una capacidad individual o una ventaja de
origen reconocida como capacidad, y si confiere merecimiento de la misma
forma que el esfuerzo.

La primera distinción sobre lo innato y lo adquirido, se evidencia que
existen diferencias genéticas entre las personas que se vinculan con
resultados como el nivel educativo alcanzado y que otorgan ventajas a
quienes las heredan, aunque su distribución sea producto del azar
(\citeproc{ref-harden_2021}{Harden, 2021}). Sin embargo, lo que se
percibe como talento no siempre es una habilidad innata. Un estudio
etnográfico sobre nadadores muestra que la excelencia no se explica por
un don especial, sino por la acumulación de prácticas y hábitos
cotidianos, de modo que el talento no sería la causa del éxito, sino la
forma de nombrarlo una vez que ya ocurrió
(\citeproc{ref-chambliss_mundanity_1989}{Chambliss, 1989}). Además, no
está claro qué es una habilidad ni cómo identificar las que requiere
cada tarea, salvo en ámbitos como el arte o el deporte
(\citeproc{ref-lampert_meritocratic_2013}{Lampert, 2013, p. 53}). Por lo
tanto, no es preciso discutir si el talento existe, sino qué se reconoce
como talento y cómo se usa para explicar posiciones y recompensas.

La segunda distinción es entre el talento como capacidad individual y
como ventaja de origen reconocida como capacidad. Según la ideología del
don, el éxito escolar de los estudiantes de origen privilegiado se
percibe como resultado de dotes naturales y no de una herencia cultural
transmitida por la familia
(\citeproc{ref-bourdieu_2009_herederos}{Bourdieu \& Passeron, 2009}).
Esa herencia es el capital cultural incorporado, que se adquiere en la
familia de origen y se expresa en disposiciones, es decir, en formas de
hablar, juzgar, percibir y actuar
(\citeproc{ref-bourdieu_distincion_2006}{Bourdieu, 2006};
\citeproc{ref-grau-grau_ilusion_2022}{Grau-Grau, 2022, p. 4}). El
sistema escolar, bajo una apariencia de neutralidad, valora estas
disposiciones como si fueran habilidades personales, y así transforma un
privilegio heredado en mérito propio
(\citeproc{ref-bourdieu_reproduccion_1998}{Bourdieu \& Passeron, 1998};
\citeproc{ref-grau-grau_ilusion_2022}{Grau-Grau, 2022, pp. 4-5}). Algo
similar ocurre en el trabajo, ya que en las ocupaciones de élite
británicas las formas de presentarse de las clases privilegiadas suelen
reconocerse como talento (\citeproc{ref-friedman_class_2019b}{Friedman
\& Laurison, 2019}). En Chile, las autoridades de universidades de élite
asumen que las credenciales y los puntajes de admisión reflejan el
talento de sus estudiantes, aun cuando esos programas son socialmente
homogéneos (\citeproc{ref-quaresma_justify_2024}{Quaresma, Villalobos,
et~al., 2024, pp. 9, 13}). Desde esta perspectiva, lo que se reconoce
como talento puede ser una ventaja de origen que se naturaliza.

La tercera distinción es si el talento confiere merecimiento de la misma
forma que el esfuerzo. Como se señaló, si el mérito es una acción libre
y esforzada, el talento solo es meritorio cuando se ejerce
(\citeproc{ref-clavero_2023}{Clavero García, 2023}). En Estados Unidos,
las personas perciben el talento como menos controlable que el esfuerzo,
pero mucho más que la suerte, y consideran que los ingresos derivados
del talento son casi tan justos como los del esfuerzo, porque a quien
usa su talento se le considera dueño de sus resultados
(\citeproc{ref-bartling_talent_2026}{Bartling et~al., 2026, p. 6}). Aun
así, como el talento se entiende como una cualidad estable que escapa al
control de las personas, explicar un resultado por el talento resulta
más fatalista que culpabilizador
(\citeproc{ref-spruyt_talent_esfuerzo_social_background_2015}{Spruyt,
2015, pp. 96-97}). Atribuir el éxito a un talento innato puede presentar
la desigualdad como inevitable
(\citeproc{ref-friedman_meaning_2024}{Friedman et~al., 2024}), y
atribuir una posición baja a la falta de talento la presenta como un
defecto natural sobre el cual la persona tiene poco control
(\citeproc{ref-sandel_tirania_2020}{Sandel, 2020};
\citeproc{ref-young_1964}{Young, 1964}). En Chile, a diferencia de lo
que ocurre con el esfuerzo, percibir que se recompensa el talento no se
asocia, o se asocia negativamente, con legitimar que el acceso a la
salud, la educación y las pensiones dependa del ingreso
(\citeproc{ref-castillo_perceptions_2025}{Castillo et~al., 2025},
\citeproc{ref-castillo_justification_desigualdades_pensiones_clases_2026}{2026}).

Las personas también recurren al talento para justificar su propia
posición, en especial quienes ocupan posiciones de élite. Entre las
élites británicas predomina una justificación basada en la aptitud
natural, la inteligencia y el buen juicio innato, mientras que las
élites danesas destacan el trabajo duro, diferencia que los autores
vinculan con el papel de las escuelas privadas de élite británicas en
cultivar la idea de un talento excepcional
(\citeproc{ref-friedman_meaning_2024}{Friedman et~al., 2024, p. 861}).
Además, entre las pocas personas de la élite británica que eran mujeres,
no blancas o de origen de clase trabajadora, se destacaba más el trabajo
duro que el talento como explicación de su éxito
(\citeproc{ref-friedman_meaning_2024}{Friedman et~al., 2024, p. 872}).
De manera similar, estudiantes de Oxford reconocen que otros tienen
menos acceso a las habilidades que se cultivan en los colegios privados,
pero atribuyen su propia admisión a una inteligencia que consideran
inherente (\citeproc{ref-warikoo_legitimating_2014}{Warikoo \& Fuhr,
2014, p. 712}). En América Latina, las élites económicas chilenas
definen el talento como habilidades empresariales y de liderazgo, más
que como inteligencia, y consideran que está disponible para todos, ya
que se trata de habilidades que cualquiera podría identificar y
reproducir (\citeproc{ref-atria_et_al_elites_meritocracia_2020}{Atria
et~al., 2020, pp. 10-11}). En Lima, en cambio, profesionales que
ascendieron socialmente reivindican su talento para distinguirse de sus
pares privilegiados, estableciendo una frontera simbólica entre el
privilegio y el talento (\citeproc{ref-renteria_telling_2026}{Rentería,
2026, pp. 1342, 1354}).

En conjunto, estos relatos muestran que el criterio que se enfatiza para
justificar la propia posición depende tanto de la posición que se ocupa
como de la socialización en instituciones como las escuelas de élite o
el mundo empresarial
(\citeproc{ref-atria_et_al_elites_meritocracia_2020}{Atria et~al.,
2020}; \citeproc{ref-friedman_meaning_2024}{Friedman et~al., 2024};
\citeproc{ref-warikoo_legitimating_2014}{Warikoo \& Fuhr, 2014}).
Además, como muestra el caso de las élites chilenas, lo que se entiende
por talento varía, por lo que al evaluar si la sociedad lo recompensa
las personas pueden estar pensando en cualidades distintas, lo que
constituye una limitación para interpretar esta percepción. Como se
trata de relatos de personas que ocupan objetivamente posiciones
aventajadas, se retoman en la sección sobre nivel socioeconómico, y su
relación con la forma en que las personas se ubican a sí mismas se
discute en la sección sobre estatus social subjetivo.

\paragraph{Esfuerzo}\label{esfuerzo-1}

El rasgo que distingue al esfuerzo es que se percibe bajo el control de
las personas. La teoría de la atribución no mide la percepción de que la
sociedad recompensa el esfuerzo, pero permite entender qué implica
explicar un resultado por él. Las personas buscan las causas de sus
éxitos y fracasos para comprender su entorno y orientar sus acciones, y
suelen atribuir el éxito a una alta capacidad y al trabajo duro, y el
fracaso a una baja capacidad y a la falta de esfuerzo
(\citeproc{ref-weiner_attributional_1985}{Weiner, 1985}). Estas causas
se clasifican según su locus, su estabilidad y su controlabilidad. El
esfuerzo y la capacidad son causas internas, pero solo el esfuerzo se
considera controlable, y es esta controlabilidad la que lleva a juzgar a
una persona como responsable de sus resultados, ya sea al evaluar a
otros o a uno mismo (\citeproc{ref-davis_learning_2017}{M. L. Davis
et~al., 2017}; \citeproc{ref-schrader_school_2015}{Schrader \& Helmke,
2015}). Así, el esfuerzo depende en mayor medida de la voluntad y la
motivación de las personas
(\citeproc{ref-franetovic_concepciones_2017a}{Franetovic, 2017, p. 5}),
y el control es uno de los criterios con que se juzga lo que otros
merecen (\citeproc{ref-meuleman_desert_carin_2020}{Meuleman et~al.,
2020}).

Kluegel \& Smith (\citeproc{ref-kluegel_beliefs_1981}{1981}) llevaron
esta teoría al estudio sociológico de las creencias sobre la desigualdad
y la combinaron con enfoques culturales, de privación relativa y de
acción colectiva para explicar cómo se atribuyen la pobreza y la
riqueza. En ese marco discuten la tesis de la ideología dominante (Huber
y Form, 1973, citado en Kluegel \& Smith
(\citeproc{ref-kluegel_beliefs_1981}{1981})). Según esta tesis, las
oportunidades están disponibles para todos, la posición de cada persona
depende de su esfuerzo y sus capacidades, y por eso la distribución de
las recompensas resulta justa. Sin embargo, la mayoría de las personas
no sostiene solo explicaciones individuales, sino que las combina con
explicaciones estructurales, en lo que se ha llamado conciencia dual,
integrando explicaciones
(\citeproc{ref-reynolds_perceptions_2014}{Reynolds \& Xian, 2014};
\citeproc{ref-zhu_meritocratic_2025}{Zhu, 2025}). Aun así, el orden
meritocrático moraliza el éxito a través del esfuerzo, y este se vuelve
el fundamento para considerar merecidas las recompensas
(\citeproc{ref-sandel_tirania_2020}{Sandel, 2020, pp. 41-60}).

Quien se esfuerza merece ser recompensado por su dedicación, pero el
esfuerzo también se usa para explicar los fracasos, con componentes de
culpa o desesperanza, porque supone que la situación estaba bajo el
control de la persona (\citeproc{ref-lampert_meritocratic_2013}{Lampert,
2013};
\citeproc{ref-spruyt_talent_esfuerzo_social_background_2015}{Spruyt,
2015, p. 96}). Un fracaso atribuido a la falta de esfuerzo genera enojo
o rechazo en otros y culpa en uno mismo, mientras que uno atribuido a la
falta de capacidad genera lástima y vergüenza
(\citeproc{ref-weiner_attributional_1985}{Weiner, 1985, p. 561}). Esta
distinción entre culpa y destino también organiza las explicaciones de
la pobreza. La explicación de culpa individual la atribuye a fallas
morales o conductuales de los propios pobres, como la pereza, la falta
de ahorro o las malas costumbres, mientras que la de destino individual
la atribuye a la mala suerte o a infortunios que están fuera del control
propio (\citeproc{ref-oorschot_blame_2000}{Oorschot \& Halman, 2000, pp.
5-7}). Así, una posición baja atribuida a la falta de esfuerzo se
convierte en una responsabilidad individual
(\citeproc{ref-sandel_tirania_2020}{Sandel, 2020}).

El esfuerzo, sin embargo, no rinde lo mismo para todos, ya que sus
límites dependen del origen social. La pequeña burguesía se caracteriza
por una buena voluntad cultural, es decir, por el esfuerzo de apropiarse
de la cultura legítima, mientras que las clases dominantes valoran la
soltura y la naturalidad, que hacen aparecer como innato lo que fue
adquirido en la familia
(\citeproc{ref-bourdieu_distincion_2006}{Bourdieu, 2006}). Del mismo
modo, en contextos escolares se valora el esfuerzo, pero lograr buenos
resultados sin esfuerzo se celebra como señal de una inteligencia
auténtica (\citeproc{ref-jackson_smart_2015}{Jackson \& Nyström, 2015,
pp. 393-394}). El esfuerzo que realizan las familias de clase
trabajadora para compensar su menor capital cultural no siempre tiene
premio (\citeproc{ref-grau-grau_ilusion_2022}{Grau-Grau, 2022, p. 5}),
mientras que quienes heredaron mayor capital cultural declaran haberse
esforzado más que el resto para obtener su posición
(\citeproc{ref-grau-grau_ilusion_2022}{Grau-Grau, 2022, pp. 15-16}).

En Chile, las trayectorias están marcadas por relatos de esfuerzo, que
operan como principio justificador, como fuente de autovaloración y como
una expectativa concreta desde la cual se evalúa la justicia de la
sociedad y el reconocimiento
(\citeproc{ref-araujo_y_martuccelli_soc-chilena_2012}{Araujo \&
Martuccelli, 2012}). En el mundo del trabajo, el esfuerzo se ha
convertido en el núcleo ético de un mercado laboral individualizado, al
punto de que ``el esfuerzo y no el mérito es lo que reclama más
compensación'' (\citeproc{ref-sir_esfuerzo_2017}{Sir, 2017, p. 311}).
Este esfuerzo excede al mérito, ya que incluye la disponibilidad, la
personalidad y las relaciones personales, en un contexto de
desconocimiento generalizado sobre los criterios de evaluación y de
desprotección laboral, donde se vuelve un requisito para ser reconocido
(\citeproc{ref-sir_esfuerzo_2017}{Sir, 2017, pp. 313, 319-320}).

Percibir que la sociedad recompensa el esfuerzo se alinea con los
relatos sobre la responsabilidad individual y el merecimiento moral
(\citeproc{ref-castillo_perceptions_2025}{Castillo et~al., 2025}). En
Chile, esta es la percepción que más se asocia con la legitimación de la
desigualdad, ya que se asocia con aceptar que el acceso a la salud, la
educación y las pensiones dependa del ingreso
(\citeproc{ref-castillo_perceptions_2025}{Castillo et~al., 2025},
\citeproc{ref-castillo_justification_desigualdades_pensiones_clases_2026}{2026}).
Por ser el criterio que permite presentar una posición como obra propia,
el esfuerzo es el que se espera que se asocie más con la forma en que
las personas se ubican y con la trayectoria que recorrieron.

\subsection{Factores asociados a la percepción de que se recompensa el
esfuerzo y el
talento}\label{factores-asociados-a-la-percepciuxf3n-de-que-se-recompensa-el-esfuerzo-y-el-talento}

\subsubsection{Nivel socioeconómico}\label{nivel-socioeconuxf3mico-1}

El nivel socioeconómico refiere a la posición de una persona según su
ingreso, su educación y su ocupación
(\citeproc{ref-li_community_2021}{Li \& Hu, 2021, p. 478}). En este
estudio se aborda con el ingreso del hogar, y la educación se trata en
una sección aparte. El ingreso permite identificar a quienes ocupan las
posiciones más altas, algo que la autoubicación no permite, porque en
Chile tiende hacia el centro
(\citeproc{ref-castillo_todos_clase_media_2013}{Castillo et~al., 2013}).

En Chile, el 10\% de mayores ingresos percibe más meritocracia que el
resto (\citeproc{ref-iturra-sanhueza_social_2026}{Iturra-Sanhueza
et~al., 2026}, Tablas S1 y S2). Esa medida combina el esfuerzo y el
talento. Cuando ambos se analizan por separado, el ingreso en deciles no
se asocia con ninguno de los dos
(\citeproc{ref-iturra-sanhueza_social_2026}{Iturra-Sanhueza et~al.,
2026}, Tabla S3), pero ese análisis no distingue a los hogares de
mayores ingresos del resto. Por eso, la pregunta es si las posiciones
altas se asocian de forma distinta con cada criterio.

En las posiciones aventajadas, las ventajas de origen tienden a
reconocerse como capacidades naturales. La soltura que se aprende en
instituciones de élite aparece como una cualidad natural
(\citeproc{ref-khan_saying_2013}{Khan \& Jerolmack, 2013, p. 16}), y las
formas de presentarse de las clases privilegiadas se reconocen como
talento en el trabajo (\citeproc{ref-friedman_class_2019b}{Friedman \&
Laurison, 2019}). De modo similar, los estudiantes de clase media y de
élite son vistos, y se ven a sí mismos, como naturalmente inteligentes,
mientras que los de clase trabajadora deben demostrar su esfuerzo
(\citeproc{ref-jackson_smart_2015}{Jackson \& Nyström, 2015, p. 402}).
Quienes ocupan posiciones aventajadas tendrían así más experiencias en
que se les reconoce por lo que tienen y no solo por lo que hacen.

La evidencia chilena también vincula las posiciones aventajadas con el
talento, donde la élite económica ancla el mérito en el talento
(\citeproc{ref-atria_et_al_elites_meritocracia_2020}{Atria et~al.,
2020}), mientras que las clases medias y bajas lo equiparan con el
esfuerzo (\citeproc{ref-araujo_y_martuccelli_soc-chilena_2012}{Araujo \&
Martuccelli, 2012};
\citeproc{ref-atria_et_al_elites_meritocracia_2020}{Atria et~al.,
2020}). Los hogares de altos ingresos no son la élite económica, que se
define por ocupar posiciones de poder en grandes empresas
(\citeproc{ref-atria_et_al_elites_meritocracia_2020}{Atria et~al.,
2020}). Pero la asociación entre mérito y talento no se limita a ella,
ya que aparece también entre directivos de universidades de élite
(\citeproc{ref-quaresma_merit_u_elites_2024}{Quaresma, Madrid, et~al.,
2024}) y, como se mencionó, entre estudiantes de clase media
(\citeproc{ref-jackson_smart_2015}{Jackson \& Nyström, 2015, p. 402}).

Existe una interpretación contraria, según la cual quienes ocupan las
posiciones más altas podrían narrar sus resultados a partir del esfuerzo
y las credenciales, y el talento serviría más a la clase intermedia
(\citeproc{ref-castillo_justification_desigualdades_pensiones_clases_2026}{Castillo
et~al., 2026, pp. 12-13}). Esa interpretación, sin embargo, se refiere a
la justificación de la desigualdad en las pensiones y no a la percepción
de que se recompensa cada criterio. Además, este argumento no supone que
en las posiciones altas desaparezca el esfuerzo, que está presente en
los relatos de todas las posiciones, incluso en las élites
(\citeproc{ref-khan_saying_2013}{Khan \& Jerolmack, 2013}). Lo que
distingue a las posiciones altas es que sus ventajas también se
reconocen como talento.

De esta manera, en posiciones altas las ventajas tenderían a reconocerse
como talentos, y las personas infieren cómo funciona la sociedad a
partir de su experiencia (\citeproc{ref-mijs_inequality_2018}{J. Mijs,
2018}). Por esto, se espera que los hogares de altos ingresos perciban
relativamente más que se recompensa el talento. No se espera que
perciban que se recompensa más el talento que el esfuerzo, sino que la
diferencia a favor del talento sea mayor que en los hogares de ingresos
bajos.

\begin{quote}
\(H_1\): Pertenecer a un hogar de altos ingresos, en comparación con uno
de ingresos bajos, se asocia de forma más positiva con la percepción de
que en Chile se recompensa el talento que con la percepción de que se
recompensa el esfuerzo.
\end{quote}

\subsubsection{Estatus social
subjetivo}\label{estatus-social-subjetivo-1}

A diferencia del nivel socioeconómico, el estatus social subjetivo
refiere a cómo las personas se ubican a sí mismas en la jerarquía
social, lo que puede no coincidir con su posición objetiva
(\citeproc{ref-davis_status_1956}{J. A. Davis, 1956, p. 154}). Como las
personas se comparan con su entorno cercano, la mayoría se ubica en el
medio (\citeproc{ref-evans_images_1992}{Evans et~al., 1992, p. 465}). En
Chile, quienes tienen un mayor estatus objetivo tienden a subestimar su
posición (\citeproc{ref-castillo_todos_clase_media_2013}{Castillo
et~al., 2013}).

En Chile, un mayor estatus subjetivo se asocia con una mayor percepción
de meritocracia, incluso al controlar por el ingreso
(\citeproc{ref-castillo_merit_perc_pref_2019}{Castillo et~al., 2019, p.
10}; \citeproc{ref-iturra-sanhueza_social_2026}{Iturra-Sanhueza et~al.,
2026}). Esto se ha explicado porque quienes están arriba en la jerarquía
suelen ser vistos como ganadores merecidos
(\citeproc{ref-iturra-sanhueza_social_2026}{Iturra-Sanhueza et~al.,
2026, p. 9}). Aun así, con datos de 2016, la asociación con la
recompensa percibida desaparece al considerar el estatus de la familia
de origen y el esperado para los hijos
(\citeproc{ref-iturrasanhueza_estatus_2019}{Iturra Sanhueza, 2019, p.
46}).

El estatus subjetivo capta cómo las personas valoran su propia posición,
más allá de sus recursos. Sentirse merecidamente arriba requiere un
criterio, y cuál es plausible depende de esos recursos. Cuando la
autoubicación no está respaldada por ventajas objetivas, presentarla
como talento es poco plausible, porque el talento se reconoce en las
posiciones objetivamente aventajadas
(\citeproc{ref-quaresma_merit_u_elites_2024}{Quaresma, Madrid, et~al.,
2024};
\citeproc{ref-atria_et_al_elites_meritocracia_2020a}{\textbf{atria\_et\_al\_elites\_meritocracia\_2020a?}}).
El esfuerzo, en cambio, permite presentarla como obra propia en
cualquier posición. Además, en Chile el esfuerzo es el lenguaje con que
las personas se valoran a sí mismas
(\citeproc{ref-araujo_y_martuccelli_soc-chilena_2012}{Araujo \&
Martuccelli, 2012}). También es la percepción del esfuerzo, y no la del
talento, la que se asocia de forma consistente con justificar
desigualdades en salud, educación y pensiones
(\citeproc{ref-castillo_perceptions_2025}{Castillo et~al., 2025},
\citeproc{ref-castillo_justification_desigualdades_pensiones_clases_2026}{2026}).

La evidencia sobre se evalúan la propia posición es la siguiente. En
Chile, quienes se ubican más arriba consideran más justo su propio
ingreso, y esa diferencia aumenta cuando lo evalúan según su esfuerzo
por educarse (\citeproc{ref-mac_baroz_merit_des_2024}{Mac-Clure et~al.,
2024, pp. 16-17}). El esfuerzo está disponible en posiciones
aventajadas, en estudiantes de la élite estadounidense explican su lugar
por su trabajo duro (\citeproc{ref-khan_saying_2013}{Khan \& Jerolmack,
2013, p. 12}), y estudiantes de universidades de élite chilenas explican
su ingreso por la fuerza de voluntad y el estudio esforzado y constante
(\citeproc{ref-villalobos_percepciones_2022}{Villalobos et~al., 2022}).

Un mayor estatus subjetivo también podría asociarse más con el talento.
En las élites chilenas el mérito se asocia al talento
(\citeproc{ref-atria_et_al_elites_meritocracia_2020}{Atria et~al.,
2020}; \citeproc{ref-quaresma_merit_u_elites_2024}{Quaresma, Madrid,
et~al., 2024}), y en Estados Unidos quienes se ubican más arriba creen
en mayor medida que la clase es algo inherente y biológicamente
determinado (\citeproc{ref-kraus_social_2013}{Kraus \& Keltner, 2013,
pp. 252, 254}). Es decir, la autoubicación también puede presentar la
posición como natural. La evidencia chilena sugiere que esto ocurre
sobre todo cuando la autoubicación está respaldada por una posición
objetivamente alta. La autoubicación solo marca diferencias en la
percepción de meritocracia sobre la mediana del ingreso
(\citeproc{ref-iturra-sanhueza_social_2026}{Iturra-Sanhueza et~al.,
2026, p. 6}). Cuando ambas percepciones se analizan por separado, el
estatus subjetivo se asocia en promedio más con el esfuerzo
(\citeproc{ref-iturra-sanhueza_social_2026}{Iturra-Sanhueza et~al.,
2026}, Tabla S3). A partir de esa tabla, en los ingresos bajos el
estatus subjetivo casi no se asocia con el esfuerzo y tiende a asociarse
negativamente con el talento. En los ingresos altos se asocia
positivamente con ambos, y la diferencia casi desaparece. Estas
diferencias no se han probado formalmente.

De esta manera, la autoubicación se legitima sobre todo mediante el
esfuerzo, porque es el criterio disponible en cualquier posición. Pero
cuando está respaldada por ingresos altos, el talento aumenta su
relevancia. Por ello se espera que el estatus subjetivo se asocie más
con el esfuerzo que con el talento, y que esta ventaja disminuya a
medida que aumenta el ingreso.

\begin{quote}
\(H_{2a}\): El estatus social subjetivo se asocia de forma más positiva
con la percepción de que en Chile se recompensa el esfuerzo que con la
percepción de que se recompensa el talento.

\(H_{2b}\): La ventaja del esfuerzo sobre el talento en la asociación
con el estatus social subjetivo disminuye a medida que aumenta el
ingreso del hogar.
\end{quote}

\subsubsection{Movilidad social
subjetiva}\label{movilidad-social-subjetiva-1}

La movilidad social subjetiva refiere a cómo las personas evalúan su
posición actual en comparación con la de la familia en la que crecieron.
Se usa la medida subjetiva y no la objetiva porque la movilidad incide
en las creencias en la medida en que las personas la experimentan como
real (\citeproc{ref-mijs_belief_2022}{Mijs et~al., 2022, pp. 132, 138}).

Quienes ascienden tienen motivos para presentar su éxito como merecido
(\citeproc{ref-mijs_belief_2022}{Mijs et~al., 2022, p. 133}). Como no
pueden atribuir su nueva posición a ventajas heredadas, tienden a
explicarla por lo que hicieron, y el esfuerzo es el criterio que permite
presentarla como obra propia
(\citeproc{ref-franetovic_concepciones_2017a}{Franetovic, 2017, p. 5}).
Además, las personas tienden a generalizar su experiencia de movilidad
al conjunto de la sociedad
(\citeproc{ref-traini_et_al_mov_social_percepciones_meritocracia_2025}{Traini
et~al., 2025, p. 4}). Por eso, haber ascendido por el propio esfuerzo
debería llevar a percibir que en la sociedad se recompensa el esfuerzo.

La relación con una misma posición también cambia según el origen de las
personas. Quienes la heredan la viven con soltura y son reconocidos como
naturalmente capaces, mientras que quienes ascienden deben apropiársela
mediante el esfuerzo (\citeproc{ref-bourdieu_distincion_2006}{Bourdieu,
2006}; \citeproc{ref-jackson_smart_2015}{Jackson \& Nyström, 2015, p.
402}). Por eso, el trabajo duro ocupa un lugar central en los relatos de
quienes ascienden. En las élites de Reino Unido y Dinamarca, quienes
ascendieron lo destacan más que el resto
(\citeproc{ref-friedman_meaning_2024}{Friedman et~al., 2024, p. 875}).
También estudiantes chinos de clase trabajadora que llegaron a
universidades de élite ponen el trabajo duro al centro de su relato, y
sostienen que siempre da frutos (\citeproc{ref-jin_meritocracy_2020}{Jin
\& Ball, 2020}). El talento también aparece en estos relatos
(\citeproc{ref-renteria_telling_2026}{Rentería, 2026, p. 1353}), pero
también se atribuye a quienes heredaron su posición, mientras que el
esfuerzo es más característico del relato del ascenso. También, se
vincula el ascenso con el trabajo duro. En Chile y otros 26 países,
quienes perciben haber ascendido atribuyen más importancia al trabajo
duro que quienes se mantuvieron en la misma posición
(\citeproc{ref-mijs_belief_2022}{Mijs et~al., 2022, p. 137};
\citeproc{ref-traini_et_al_mov_social_percepciones_meritocracia_2025}{Traini
et~al., 2025, pp. 11-12}).

Sin embargo, estos estudios miden la importancia atribuida al trabajo
duro y no la percepción de que se recompensa, dos dimensiones distintas
en Chile (\citeproc{ref-iturrasanhueza_estatus_2019}{Iturra Sanhueza,
2019, p. 35}), y tampoco separan el talento. Además, en Chile quienes
tuvieron una movilidad ascendente de largo alcance reconocen los
factores externos que la hicieron posible
(\citeproc{ref-fercovic_disentanglin_2022}{Fercovic, 2022}).

La movilidad subjetiva está ligada al estatus subjetivo, porque compara
la posición actual con la de origen, y el estatus de origen es el que
más pesa en cómo las personas se ubican
(\citeproc{ref-iturra_status_subj_2019}{Iturra-Sanhueza \& Riffo, 2019,
p. 19}). Por eso, para distinguir la trayectoria de la posición
alcanzada, la comparación relevante es entre personas que ocupan una
misma posición, pero llegaron a ella por caminos distintos. Unas
ascendieron desde un origen más bajo y otras la heredaron, porque ya
provenían de ella.

Este argumento debe conciliarse con la evidencia chilena sobre el
estatus de origen. Entre personas que se ubican en el mismo nivel,
quienes provienen de familias de mayor estatus perciben en mayor medida
que en Chile se recompensan el esfuerzo y la inteligencia
(\citeproc{ref-iturrasanhueza_estatus_2019}{Iturra Sanhueza, 2019, p.
46}). Es decir, quienes heredaron su posición perciben más recompensa
que quienes ascendieron hasta ella. Sin embargo, como esa medida combina
ambos criterios, esa ventaja podría concentrarse en el talento. De esta
manera, entre personas que ocupan una misma posición, quienes llegaron a
ella ascendiendo tienen motivos y experiencias para asociarla al
esfuerzo. Quienes ya provenían de ella tienen más experiencias en que se
les reconoce como naturalmente capaces. Por ello se espera que
experimentar movilidad ascendente, en comparación con no experimentarla,
se asocie relativamente más con la percepción de que se recompensa el
esfuerzo. La movilidad descendente, cuya evidencia es menos consistente
(\citeproc{ref-mijs_belief_2022}{Mijs et~al., 2022, p. 137};
\citeproc{ref-traini_et_al_mov_social_percepciones_meritocracia_2025}{Traini
et~al., 2025, p. 12}). Así, se considera sin una hipótesis.

\begin{quote}
\(H_3\): Experimentar movilidad social subjetiva ascendente, en
comparación con no experimentar movilidad, se asocia de forma más
positiva con la percepción de que en Chile se recompensa el esfuerzo que
con la percepción de que se recompensa el talento.
\end{quote}

\subsubsection{Nivel educacional}\label{nivel-educacional-1}

La educación ocupa un lugar central en la meritocracia. Según la teoría
de la meritocracia basada en la educación, la selección social depende
cada vez más del logro individual y de las credenciales educativas
(\citeproc{ref-goldthorpe_myth_2003}{Goldthorpe, 2003, p. 234}).

En la educación y las percepciones meritocráticas hay dos hipótesis
opuestas que pueden operar a la vez, la socialización y la iluminación.
En la primera, la escuela transmite que el éxito depende del mérito
individual, por lo que quienes avanzan más en el sistema educativo
adhieren más a la meritocracia
(\citeproc{ref-bourdieu_reproduccion_1998}{Bourdieu \& Passeron, 1998};
\citeproc{ref-lampert_meritocratic_2013}{Lampert, 2013}). Y la segunda,
la educación aumentaría la sensibilidad frente a las desigualdades, por
lo que los más educados son más críticos respecto de cómo se distribuyen
las recompensas (\citeproc{ref-bucca_merit_2016}{Bucca, 2016};
\citeproc{ref-duru-bellat_variations_2012}{Duru-Bellat \& Tenret, 2012};
\citeproc{ref-kluegel_beliefs_1981}{Kluegel \& Smith, 1981}). En Chile,
un mayor nivel educacional se asocia con una menor percepción de
meritocracia (\citeproc{ref-castillo_merit_perc_pref_2019}{Castillo
et~al., 2019, p. 9};
\citeproc{ref-iturra-sanhueza_social_2026}{Iturra-Sanhueza et~al.,
2026}). Las personas más educadas atribuyen más importancia a la
ambición, la educación y el trabajo duro para salir adelante, como
menciona la socialización
(\citeproc{ref-iturrasanhueza_estatus_2019}{Iturra Sanhueza, 2019, pp.
41-42}). Pero, y siguiendo la segunda hipótesis, perciben en menor
medida que se recompensan el esfuerzo y la inteligencia
(\citeproc{ref-iturrasanhueza_estatus_2019}{Iturra Sanhueza, 2019, p.
45}). En la percepción de recompensa, entonces, predomina la
iluminación. La pregunta es si lo hace con la misma fuerza para el
esfuerzo y para el talento.

La iluminación afectaría ambas percepciones, porque los más educados
reconocen las barreras que impiden que el mérito se traduzca en
recompensas y la socialización debería pesar más en el talento. La
escuela traduce las ventajas de origen en capacidades individuales
(\citeproc{ref-accominotti_making_2026}{Accominotti \& Sauder, 2026}) y
valora como mérito propio lo que es un habitus de clase
(\citeproc{ref-grau-grau_ilusion_2022}{Grau-Grau, 2022, p. 4}). Así,
quienes avanzan en el sistema educativo reciben credenciales que
certifican su capacidad, y tienen más experiencias en que su talento es
reconocido. Por eso, en el caso del talento, la socialización
compensaría en parte la crítica que produce la iluminación.

La evidencia para ambos criterios es limitada. Cuando se mide la
percepción de que se recompensa cada talento y esfuerzo, la educación se
asocia negativamente con ambos, con magnitudes similares, aunque algo
menos para el talento en tres de cuatro niveles educacionales
(\citeproc{ref-iturra-sanhueza_social_2026}{Iturra-Sanhueza et~al.,
2026}, Tabla S3). En cambio, en Flandes, son los más educados quienes
recurren menos al talento para explicar éxito escolar
(\citeproc{ref-spruyt_talent_esfuerzo_social_background_2015}{Spruyt,
2015}). En Chile, la educación universitaria de los padres se asocia con
una menor importancia del talento para obtener buenas notas
(\citeproc{ref-retamales_rol_2024}{Retamales, 2024, pp. 44-45}). Pero en
ambos países se estudian las explicaciones de éxito escolar, no las
percepciones de recompensa según los dos criterios.

Entonces, se plantea que la educación haría más crítica ambas
percepciones, pero que la certificación escolar de las capacidades
atenuaría esa crítica en el caso del talento. Por esto, se espera que la
educación se asocie negativamente con ambas percepciones, pero en menor
medida con la de talento.

\begin{quote}
\(H_4\): La asociación negativa del nivel educacional es más débil con
la percepción de que en Chile se recompensa el talento que con la
percepción de que se recompensa el esfuerzo.
\end{quote}

\subsubsection{Sexo y edad}\label{sexo-y-edad}

Respecto del sexo, no existe consenso sobre su relación con la
percepción de meritocracia. Los hombres creen en mayor medida que se
recompensan los logros
(\citeproc{ref-duru-bellat_variations_2012}{Duru-Bellat \& Tenret, 2012,
p. 233}), pero en Chile la evidencia no es consistente, con datos de
2016, las mujeres perciben menos meritocracia
(\citeproc{ref-plaza_networked_2026}{Plaza et~al., 2026, p. 35}) o no se
diferencian de los hombres
(\citeproc{ref-iturrasanhueza_estatus_2019}{Iturra Sanhueza, 2019, p.
70}). Con datos de 2018, perciben más que se recompensan tanto el
esfuerzo como el talento
(\citeproc{ref-iturra-sanhueza_social_2026}{Iturra-Sanhueza et~al.,
2026}, Tabla S3). De todas maneras, en la escuela, los docentes tienden
a atribuir el éxito de los niños a su capacidad y el de las niñas a su
esfuerzo
(\citeproc{ref-fennema_teachers_1990}{\textbf{fennema\_teachers\_1990?}};
\citeproc{ref-espinoza_attributional_2014}{\textbf{espinoza\_attributional\_2014?}}).
El logro sin esfuerzo aparente, también se asocia con la masculinidad
(\citeproc{ref-jackson_smart_2015}{Jackson \& Nyström, 2015, p. 402};
\citeproc{ref-jackson_dempster_2009}{\textbf{jackson\_dempster\_2009?}}).
Entre estudiantes chilenos, los hombres dan más importancia al talento y
menos al esfuerzo que las mujeres para obtener buenas notas
(\citeproc{ref-retamales_rol_2024}{Retamales, 2024, pp. 40, 46}). Si las
personas infieren cómo funciona la sociedad a partir de su experiencia,
los hombres podrían percibir relativamente más que se recompensa el
talento. Sin embargo, entre adultos de Flandes, hombres y mujeres no se
diferencian en cómo explican el éxito escolar
(\citeproc{ref-spruyt_talent_esfuerzo_social_background_2015}{Spruyt,
2015, p. 106}).

Respecto de la edad, las personas mayores creen en mayor medida que se
recompensan los logros
(\citeproc{ref-duru-bellat_variations_2012}{Duru-Bellat \& Tenret, 2012,
p. 233}). En Chile, la edad se asocia levemente con una mayor percepción
de meritocracia (\citeproc{ref-iturrasanhueza_estatus_2019}{Iturra
Sanhueza, 2019, p. 70}; \citeproc{ref-plaza_networked_2026}{Plaza
et~al., 2026, p. 36}), o no se asocia con ninguna de las dos
percepciones (\citeproc{ref-iturra-sanhueza_social_2026}{Iturra-Sanhueza
et~al., 2026}). Sin embargo, podría marcar una diferencia entre ambos
criterios. En Flandes, las personas mayores explican el éxito escolar
mucho más por el talento, pero casi no se diferencian en el esfuerzo
(\citeproc{ref-spruyt_talent_esfuerzo_social_background_2015}{Spruyt,
2015}). Spruyt
(\citeproc{ref-spruyt_talent_esfuerzo_social_background_2015}{2015})
sugiere que estas generaciones fueron a la escuela cuando la educación
recién se masificaba, y luego vieron ampliarse las oportunidades. Por
eso creerían que hoy el éxito depende sobre todo del talento
(\citeproc{ref-spruyt_talent_esfuerzo_social_background_2015}{Spruyt,
2015, p. 105}). Así, las personas mayores podrían percibir relativamente
más que se recompensa el talento.

Dado que la evidencia no es concluyente, el sexo y la edad se incluyen
en el análisis sin una hipótesis específica.

\bookmarksetup{startatroot}

\chapter{Metodología}\label{metodologuxeda-1}

\section{Datos}\label{datos-1}

En el presente estudio, se utilizarán los datos de la encuesta
``Economía Moral de la Meritocracia y Preferencias Redistributivas''
parte del proyecto Fondecyt 1160921, llevada a cabo en abril 2016 -
marzo 2020. Los análisis se basan en datos de la primera ola,
recopilados entre el 25 de noviembre de 2019 y el 8 de enero de 2020. La
muestra no es probabilística. Se construyó con cuotas de sexo, edad y
nivel educacional, basadas en la encuesta CEP de 2019
(\citeproc{ref-castillo_et_al_multidimensional_merit_beliefs_2023}{Castillo
et~al., 2023, p. 461}). Esta encuesta tiene una muestra de N = 2.236,
incluyendo residentes de 54 comunas. La base no incluye ponderadores. Y,
contempla un cuestionario de 37 preguntas aplicado de manera online
mediante Qualtrics. Los encuestados responden preguntas sobre temas como
tolerancia a la desigualdad, importancia para surgir en la vida,
igualitarismo, percepciones y preferencias en meritocracia, entre otros.
Al ser aplicado de manera online no probabilístico, excluye a personas
que no usan internet, y subrepresenta a las personas de menor nivel
educacional. Por esto, los resultados no se generalizan a la población
chilena y se interpretan como asociaciones dentro de la muestra.

La encuesta tiene dos puntos fuertes en relación con la investigación de
meritocracia. Incluye percepciones y preferencias meritocráticas y junto
a este, factores no meritocráticos, lo que permite dimensionar el
concepto más allá de las creencias y también separa el concepto en sus
dos dimensiones, el talento y esfuerzo. También, se utiliza la primera
ola, ya que es la única que incluye ingresos del hogar, y se responde a
preguntas de meritocracia.

\section{Variables dependientes: Percepciones de esfuerzo y
talento}\label{variables-dependientes-percepciones-de-esfuerzo-y-talento-1}

Las variables dependientes son la percepción de que en Chile se
recompensa el esfuerzo y la percepción de que se recompensa el talento.
Cada una se mide con un ítem tipo Likert de 1 (totalmente en desacuerdo)
a 5 (totalmente de acuerdo), en que valores más altos indican una mayor
percepción de que se recompensa ese criterio.

\section{Variables independientes}\label{variables-independientes-1}

El nivel socioeconómico se mide con el ingreso líquido mensual del
hogar, en 14 tramos. Se recodificará siguiendo la estrategia de estimar
diferencias entre grupos de ingreso en lugar de un coeficiente único
(\citeproc{ref-mac_baroz_merit_des_2024}{Mac-Clure et~al., 2024, p.
13}), porque no se puede construir quintiles de ingreso per cápita al no
preguntar por el número de integrantes del hogar. Los tramos se agrupan
en cuatro categorías: ingresos bajos (hasta \$448.000, referencia),
medio-bajos (\$448.001 a \$1.000.000), medio-altos (\$1.000.001 a
\$2.000.000) y altos (más de \$2.000.000). Esta última se aproxima al
10\% de mayores ingresos que distinguen Iturra-Sanhueza et~al.
(\citeproc{ref-iturra-sanhueza_social_2026}{2026}).

El estatus social subjetivo que se mide como una escala de autoubicación
de 11 puntos donde se pregunta: `En Chile, hay grupos que tienden a
ubicarse en los niveles más altos y grupos que tienden a ubicarse en los
niveles más bajos de la sociedad. Usando la escala presentada, donde 0
es el nivel más bajo y 10 el nivel más alto, ¿dónde se ubicaría usted en
la sociedad chilena?'. En la cual valores más altos significan mayor
estatus social percibido.

La movilidad social subjetiva se construye con la variable de estatus
subjetivo propio y el de la familia de origen, el ítem menciona: `\,``Y
pensando en la familia en la que usted creció, ¿dónde se ubicarían ellos
en esta escala?'', de 0 a 10'. Siguiendo a Mijs et~al.
(\citeproc{ref-mijs_belief_2022}{2022, p. 134}) y Kwon \& Erola
(\citeproc{ref-kwon_access_2026}{2026, p. 1008}), ambas escalas se
agruparán en bajo (0 a 3), medio (4 a 6) y alto (7 a 10) y se cruzan.
Quienes quedan en un nivel superior al de su familia se clasifican como
ascendentes, quienes quedan en uno inferior como descendentes y quienes
se mantienen como inmóviles.

El nivel educacional se mide con el mayor nivel alcanzado, en 10
categorías. Se recodifica según el último nivel completado, con base en
la Clasificación Internacional Normalizada de la Educación
(\citeproc{ref-statistics_clasificacion_2013}{Instituto de Estadística
de la UNESCO, 2013}), como en Iturra-Sanhueza et~al.
(\citeproc{ref-iturra-sanhueza_social_2026}{2026}). Las categorías son
secundaria alta o menos (hasta media completa, además de la educación
técnica o universitaria incompleta) de referencia, terciaria de ciclo
corto (técnica completa) y terciaria o más (universitaria completa o
postgrado). Los niveles inferiores a la secundaria alta se agrupan con
ella por poca cantidad de casos (8,8\% en total). Se trata como variable
categórica, porque la distancia entre niveles no es necesariamente la
misma.

\section{Variables de control}\label{variables-de-control-2}

La variable de sexo es dicotómica donde (1) es hombre y (2) es mujer, se
recodifica a variable dummy, donde (0) es hombre, la cual se utilizará
como categoría de referencia y (1) es mujer. La edad se utilizará como
una variable continua. Además, se incluye la versión de la batería de
meritocracia, que se presentó en tres órdenes distintos asignados al
azar
(\citeproc{ref-castillo_et_al_multidimensional_merit_beliefs_2023}{Castillo
et~al., 2023, p. 462}).

La Tabla~\ref{tbl-descriptivos} presenta los estadísticos descriptivos
de todas las variables del estudio.

\begin{Shaded}
\begin{Highlighting}[]
\NormalTok{tabla\_desc }\OtherTok{\textless{}{-}} \FunctionTok{readRDS}\NormalTok{(}\FunctionTok{here}\NormalTok{(}\StringTok{"input"}\NormalTok{, }\StringTok{"data{-}proc"}\NormalTok{, }\StringTok{"tabla\_desc.rds"}\NormalTok{))}
\FunctionTok{kable}\NormalTok{(tabla\_desc,}
      \AttributeTok{col.names =} \FunctionTok{c}\NormalTok{(}\StringTok{"Variable"}\NormalTok{, }\StringTok{"Categoría"}\NormalTok{, }\StringTok{"n o media"}\NormalTok{, }\StringTok{"\% o DE"}\NormalTok{, }\StringTok{"Rango"}\NormalTok{),}
      \AttributeTok{align =} \FunctionTok{c}\NormalTok{(}\StringTok{"l"}\NormalTok{, }\StringTok{"l"}\NormalTok{, }\StringTok{"r"}\NormalTok{, }\StringTok{"r"}\NormalTok{, }\StringTok{"c"}\NormalTok{))}
\end{Highlighting}
\end{Shaded}

{

\begin{longtable}[]{@{}
  >{\raggedright\arraybackslash}p{(\linewidth - 8\tabcolsep) * \real{0.3978}}
  >{\raggedright\arraybackslash}p{(\linewidth - 8\tabcolsep) * \real{0.3333}}
  >{\raggedleft\arraybackslash}p{(\linewidth - 8\tabcolsep) * \real{0.1075}}
  >{\raggedleft\arraybackslash}p{(\linewidth - 8\tabcolsep) * \real{0.0753}}
  >{\centering\arraybackslash}p{(\linewidth - 8\tabcolsep) * \real{0.0860}}@{}}

\caption{\label{tbl-descriptivos}Estadísticos descriptivos de las
variables del estudio (N = 2.079)}

\tabularnewline

\toprule\noalign{}
\begin{minipage}[b]{\linewidth}\raggedright
Variable
\end{minipage} & \begin{minipage}[b]{\linewidth}\raggedright
Categoría
\end{minipage} & \begin{minipage}[b]{\linewidth}\raggedleft
n o media
\end{minipage} & \begin{minipage}[b]{\linewidth}\raggedleft
\% o DE
\end{minipage} & \begin{minipage}[b]{\linewidth}\centering
Rango
\end{minipage} \\
\midrule\noalign{}
\endhead
\bottomrule\noalign{}
\endlastfoot
Percepción de recompensa al esfuerzo & & 3.20 & 1.38 & 1 a 5 \\
Percepción de recompensa al talento & & 3.03 & 1.16 & 1 a 5 \\
Ingreso del hogar & Bajos (ref.) & 686 & 33.0 & \\
& Medio-bajos & 744 & 35.8 & \\
& Medio-altos & 419 & 20.2 & \\
& Altos & 230 & 11.1 & \\
Nivel educacional & Secundaria alta o menos (ref.) & 1409 & 67.8 & \\
& Terciaria de ciclo corto & 249 & 12.0 & \\
& Terciaria o más & 421 & 20.3 & \\
Estatus social subjetivo & & 4.37 & 1.79 & 0 a 10 \\
Movilidad social subjetiva & Ascendente & 539 & 25.9 & \\
& Descendente & 194 & 9.3 & \\
& NA & 1346 & 64.7 & \\
Sexo & Hombre (0) & 1041 & 50.1 & \\
& Mujer (1) & 1038 & 49.9 & \\
Edad & 18-24 (ref.) & 387 & 18.6 & \\
& 25-34 & 389 & 18.7 & \\
& 35-44 & 395 & 19.0 & \\
& 45-54 & 374 & 18.0 & \\
& 55 o más & 534 & 25.7 & \\
Versión de la batería & A & 693 & 33.3 & \\
& B & 700 & 33.7 & \\
& C & 686 & 33.0 & \\

\end{longtable}

}

\section{Métodos}\label{muxe9todos-1}

\textbf{Breve descripción de los procedimientos aplicados para procesar
y analizar la información.}

Se estima una regresión lineal multivariante en la que ambas
percepciones son variables dependientes simultáneas, con las mismas
variables independientes. Esto permite estimar si los factores asociados
de ambas variables dependientes afectan de igual manera. Test de Wald.
Los análisis se realizan en R.

\textbf{La movilidad, el estatus subjetivo propio y el de origen están
relacionados por construcción, por lo que no pueden separarse por
completo. Por eso, \(H_3\) se contrasta en dos especificaciones. En la
primera, sin el estatus subjetivo propio, el coeficiente de la movilidad
compara a quienes ascendieron con quienes se mantuvieron, e incluye el
hecho de que quienes ascienden tienden a ubicarse más arriba. En la
segunda, que incluye el estatus propio, el coeficiente compara a
personas que se ubican en el mismo nivel pero provienen de orígenes
distintos, es decir, refleja sobre todo la diferencia de origen. \(H_3\)
se considera respaldada si la diferencia a favor del esfuerzo aparece en
ambas especificaciones. La colinealidad entre la movilidad y el estatus
propio se revisa con el factor de inflación de la varianza.}

\bookmarksetup{startatroot}

\chapter{Resultados}\label{resultados}

\section{Analisis descriptivo}\label{analisis-descriptivo}

Resumen de las estadisticas descriptivas clave y visualizaciones de
apoyo.

\section{Modelos}\label{modelos}

Presentacion de los resultados de los modelos propuestos, con
interpretacion y conclusiones.

\bookmarksetup{startatroot}

\chapter{Discusión}\label{discusiuxf3n}

Interpreta resultados, limitaciones y aportes.

\bookmarksetup{startatroot}

\chapter{Conclusiones}\label{conclusiones}

Síntesis y proyecciones futuras.

\bookmarksetup{startatroot}

\chapter{Referencias}\label{referencias}

\protect\phantomsection\label{refs}
\begin{CSLReferences}{1}{0}
\bibitem[\citeproctext]{ref-accominotti_making_2026}
Accominotti, F., \& Sauder, M. (2026). The {Making} of {Meritocratic
Status Orders}. \emph{Annual Review of Sociology}, \emph{52}(1),
301-318. \url{https://doi.org/10.1146/annurev-soc-090225-014158}

\bibitem[\citeproctext]{ref-araujo_y_martuccelli_soc-chilena_2012}
Araujo, K., \& Martuccelli, D. (2012). \emph{{Desafíos comunes: retrato
de la sociedad chilena y sus individuos}} (1a. ed). LOM Ediciones.

\bibitem[\citeproctext]{ref-atria_et_al_elites_meritocracia_2020}
Atria, J., Castillo, J., Maldonado, L., \& Ramirez, S. (2020). Economic
{Elites}' {Attitudes Toward Meritocracy} in {Chile}: {A Moral Economy
Perspective}. \emph{American Behavioral Scientist}, \emph{64}(9),
1219-1241. \url{https://doi.org/10.1177/0002764220941214}

\bibitem[\citeproctext]{ref-bartling_talent_2026}
Bartling, B., Cappelen, A. W., Skarpeid, I. L., Sørensen, E. Ø., \&
Tungodden, B. (2026). The Talent Paradox: {Why} Is It Fair to Reward
Talent but Not Luck? \emph{European Economic Review}, \emph{185},
105268. \url{https://doi.org/10.1016/j.euroecorev.2026.105268}

\bibitem[\citeproctext]{ref-bourdieu_distincion_2006}
Bourdieu, P. (2006). \emph{{La distinción: criterios y bases sociales
del gusto}} (3. ed). Taurus.

\bibitem[\citeproctext]{ref-bourdieu_reproduccion_1998}
Bourdieu, P., \& Passeron, J.-C. (1998). \emph{La Reproducción Elementos
Para Una Teoría Del Sistema de Enseñanza} (3a. ed). La Fontamara.

\bibitem[\citeproctext]{ref-bourdieu_2009_herederos}
Bourdieu, P., \& Passeron, J.-C. (2009). \emph{Los Herederos: Los
Estudiantes y La Cultura} (M. Mayer, Trad.; 2.ª ed., p. 216). Siglo XXI
Editores.

\bibitem[\citeproctext]{ref-bucca_merit_2016}
Bucca, M. (2016). Merit and Blame in Unequal Societies: {Explaining
Latin Americans}' Beliefs about Wealth and Poverty. \emph{Research in
Social Stratification and Mobility}, \emph{44}, 98-112.
\url{https://doi.org/10.1016/j.rssm.2016.02.005}

\bibitem[\citeproctext]{ref-castillo_justification_desigualdades_pensiones_clases_2026}
Castillo, J. C., Canales-Sellés, R., Laffert, A., \& Urzúa, T. (2026).
Justification Trajectories for Pension Inequality in {Chile}
(2016--2023): The Role of Social Class and Beliefs in Meritocracy.
\emph{Front. Sociol.}, \emph{11}.
\url{https://doi.org/10.3389/fsoc.2026.1771856}

\bibitem[\citeproctext]{ref-castillo_et_al_multidimensional_merit_beliefs_2023}
Castillo, J. C., Iturra, J., Maldonado, L., Atria, J., \& Meneses, F.
(2023). A {Multidimensional Approach} for {Measuring Meritocratic
Beliefs}: {Advantages}, {Limitations} and {Alternatives} to the {ISSP
Social Inequality Survey}. \emph{International Journal of Sociology},
\emph{53}(6), 448-472.
\url{https://doi.org/10.1080/00207659.2023.2274712}

\bibitem[\citeproctext]{ref-castillo_perceptions_2025}
Castillo, J. C., Laffert, A., Carrasco, K., \& Iturra-Sanhueza, J.
(2025). Perceptions of Inequality and Meritocracy: Their Interplay in
Shaping Preferences for Market Justice in {Chile} (2016--2023).
\emph{Front. Sociol.}, \emph{10}.
\url{https://doi.org/10.3389/fsoc.2025.1634219}

\bibitem[\citeproctext]{ref-castillo_todos_clase_media_2013}
Castillo, J. C., Miranda, D., \& Cabib, I. M. (2013). Todos Somos de
Clase Media: {Sobre} El Estatus Social Subjetivo En {Chile}. \emph{Latin
American Research Review}, \emph{48}(1), 155-173.
\url{https://muse.jhu.edu/pub/101/article/504733}

\bibitem[\citeproctext]{ref-castillo_percepcion_2012}
Castillo, J. C., Miranda, D., \& Carrasco, D. (2012). Percepción de
{Desigualdad Económica} En {Chile}: {Medición}, {Diferencias} y
{Determinantes}. \emph{Psykhe (Santiago)}, \emph{21}(1), 99-114.
\url{https://doi.org/10.4067/S0718-22282012000100007}

\bibitem[\citeproctext]{ref-castillo_merit_perc_pref_2019}
Castillo, J. C., Torres, A., Atria, J., \& Maldonado, L. (2019).
{Meritocracia y desigualdad económica: Percepciones, preferencias e
implicancias}. \emph{Rev. int. sociol.}, \emph{77}(1), e117.
\url{https://doi.org/10.3989/ris.2019.77.1.17.114}

\bibitem[\citeproctext]{ref-chambliss_mundanity_1989}
Chambliss, D. F. (1989). The {Mundanity} of {Excellence}: {An
Ethnographic Report} on {Stratification} and {Olympic Swimmers}.
\emph{Sociological Theory}, \emph{7}(1), 70-86.
\url{https://doi.org/10.2307/202063}

\bibitem[\citeproctext]{ref-clavero_2023}
Clavero García, S. R. (2023). The {Idea} of {Merit}: {Delineation} and
{Challenges}. \emph{Philosophia}, \emph{51}(3), 1175-1191.
\url{https://doi.org/10.1007/s11406-023-00628-x}

\bibitem[\citeproctext]{ref-davis_status_1956}
Davis, J. A. (1956). Status {Symbols} and the {Measurement} of {Status
Perception}. \emph{Sociometry}, \emph{19}(3), 154.
\url{https://doi.org/10.2307/2785629}

\bibitem[\citeproctext]{ref-davis_est_1970}
Davis, K., \& Moore, W. E. (1970). {Algunos principios de la teor{ı́}a de
la estratificación}. En R. Bendix \& S. M. Lipset (Eds.), \emph{{La
estructura de clases}}. Editorial Tiempo Nuevo.

\bibitem[\citeproctext]{ref-davis_learning_2017}
Davis, M. L., Witcraft, S. M., Baird, S. O., \& Smits, J. A. J. (2017).
Learning {Principles} in {CBT}. En \emph{The {Science} of {Cognitive
Behavioral Therapy}} (pp. 51-76). Elsevier.
\url{https://doi.org/10.1016/B978-0-12-803457-6.00003-9}

\bibitem[\citeproctext]{ref-dawtry_why_2015}
Dawtry, R. J., Sutton, R. M., \& Sibley, C. G. (2015). Why {Wealthier
People Think People Are Wealthier}, and {Why It Matters}: {From Social
Sampling} to {Attitudes} to {Redistribution}. \emph{Psychol Sci},
\emph{26}(9), 1389-1400. \url{https://doi.org/10.1177/0956797615586560}

\bibitem[\citeproctext]{ref-dubet_nuevo_2023}
Dubet, F. (2023). \emph{{El Nuevo Régimen de Las Desigualdades
Solitarias: Qué Hacer Cuando la Injusticia Social Se Sufre Como un
Problema Individual}} (1st ed). Siglo XXI Editores.

\bibitem[\citeproctext]{ref-duru-bellat_variations_2012}
Duru-Bellat, M., \& Tenret, E. (2012). Who's for {Meritocracy}?
{Individual} and {Contextual Variations} in the {Faith}.
\emph{Comparative Education Review}, \emph{56}(2), 223-247.
\url{https://doi.org/10.1086/661290}

\bibitem[\citeproctext]{ref-ehmes_paradox_2026}
Ehmes, S., \& Gangl, M. (2026). The Paradox of Inequality That Isn't:
Rising Economic Inequality Depresses and Polarizes Citizens' Belief in
Meritocracy. \emph{Socio-Economic Review}, \emph{24}(3), 1643-1663.
\url{https://doi.org/10.1093/ser/mwag016}

\bibitem[\citeproctext]{ref-evans_images_1992}
Evans, M. D. R., Kelley, J., \& Kolosi, T. (1992). Images of {Class}:
{Public Perceptions} in {Hungary} and {Australia}. \emph{American
Sociological Review}, \emph{57}(4), 461.
\url{https://doi.org/10.2307/2096095}

\bibitem[\citeproctext]{ref-fercovic_disentanglin_2022}
Fercovic, M. (2022). Disentangling {Meritocracy Among} the {Long-Range
Upwardly Mobile}: {The Chilean Case}. \emph{Sociological Research
Online}, \emph{27}(1), 118-135.
\url{https://doi.org/10.1177/1360780420963395}

\bibitem[\citeproctext]{ref-franetovic_concepciones_2017a}
Franetovic, G. (2017). \emph{Concepciones de Mérito y Meritocracia}
{[}Working Paper{]}. Pontificia Universidad Católica de Chile.
\url{https://doi.org/10.13140/RG.2.2.14983.42401/1}

\bibitem[\citeproctext]{ref-friedman_meaning_2024}
Friedman, S., Ellersgaard, C., Reeves, A., \& Larsen, A. G. (2024). The
{Meaning} of {Merit}: {Talent} versus {Hard Work Legitimacy}. \emph{Soc
Forces}, \emph{102}(3), 861-879.
\url{https://doi.org/10.1093/sf/soad131}

\bibitem[\citeproctext]{ref-friedman_class_2019b}
Friedman, S., \& Laurison, D. (2019). The Class Pay Gap: Why It Pays to
Be Privileged. \emph{The Guardian}.
\url{https://www.theguardian.com/society/2019/feb/07/the-class-pay-gap-why-it-pays-to-be-privileged}

\bibitem[\citeproctext]{ref-friedman_deflecting_2021a}
Friedman, S., O'Brien, D., \& McDonald, I. (2021). Deflecting
{Privilege}: {Class Identity} and the {Intergenerational Self}.
\emph{Sociology}, \emph{55}(4), 716-733.
\url{https://doi.org/10.1177/0038038520982225}

\bibitem[\citeproctext]{ref-gogescu_meritocracy_2025}
Gogescu, F. (2025). Meritocracy Seen through the Eyes of Its Champions:
A Comparative Study of Educational Elites. \emph{British Journal of
Sociology of Education}, \emph{46}(1), 34-53.
\url{https://doi.org/10.1080/01425692.2024.2418097}

\bibitem[\citeproctext]{ref-goldthorpe_myth_2003}
Goldthorpe, J. (2003). The Myth of Education-based Meritocracy.
\emph{New Economy}, \emph{10}(4), 234-239.
\url{https://doi.org/10.1046/j.1468-0041.2003.00324.x}

\bibitem[\citeproctext]{ref-grau-grau_ilusion_2022}
Grau-Grau, M. (2022). La Ilusión Del Mérito. {Capital} Cultural y
Reproducción. \emph{RES}, \emph{31}(3), a124.
\url{https://doi.org/10.22325/fes/res.2022.124}

\bibitem[\citeproctext]{ref-harden_2021}
Harden, K. P. (2021). \emph{{La lotería genética}}. Deusto.

\bibitem[\citeproctext]{ref-hauhart_davis_2003}
Hauhart, R. C. (2003). The {Davis}--{Moore} Theory of Stratification:
{The} Life Course of a Socially Constructed Classic. \emph{The American
Sociologist}, \emph{34}, 5-24.
\url{https://doi.org/10.1007/s12108-003-1013-y}

\bibitem[\citeproctext]{ref-cfps_data}
Institute of Social Science Survey, Peking University. (2015).
\emph{China {Family Panel Studies} ({CFPS})}. Peking University Open
Research Data Platform. \url{https://doi.org/10.18170/DVN/45LCSO}

\bibitem[\citeproctext]{ref-statistics_clasificacion_2013}
Instituto de Estadística de la UNESCO. (2013). \emph{Clasificación
Internacional Normalizada de La Educación, {CINE} 2011}. UNESCO.
\url{https://unesdoc.unesco.org/ark:/48223/pf0000220782}

\bibitem[\citeproctext]{ref-issp_data}
ISSP Research Group. (2002). \emph{International {Social Survey
Programme}: {Social Inequality III} - {ISSP 1999International Social
Survey Programme}: {Social Inequality III} - {ISSP} 1999}. GESIS Data
Archive. \url{https://doi.org/10.4232/1.3430}

\bibitem[\citeproctext]{ref-iturrasanhueza_estatus_2019}
Iturra Sanhueza, J. C. (2019). \emph{{Estatus social subjetivo y
percepción de meritocracia: el caso de Chile}} {[}\{Tesis de
mag\textbackslash\textquotesingle ister\}, Pontificia Universidad
Católica de Chile{]}. \url{https://doi.org/10.7764/tesisUC/SOC/65873}

\bibitem[\citeproctext]{ref-iturra-sanhueza_social_2026}
Iturra-Sanhueza, J., Castillo, J. C., \& García-Castro, J. D. (2026).
Social Hierarchies and Meritocracy: Objective Status and the Moderating
Role of Subjective Social Status on Perceived Meritocracy. \emph{Front.
Sociol.}, \emph{11}, 1750965.
\url{https://doi.org/10.3389/fsoc.2026.1750965}

\bibitem[\citeproctext]{ref-iturra_status_subj_2019}
Iturra-Sanhueza, J., \& Riffo, D. M. (2019). \emph{Subjective Social
Status in Three {Latin American} Countries: {The} Case of {Argentina},
{Chile} and {Venezuela}.} SocArXiv.
\url{https://doi.org/10.31235/osf.io/yp3wm}

\bibitem[\citeproctext]{ref-jackson_smart_2015}
Jackson, C., \& Nyström, A.-S. (2015). {«{Smart} Students Get Perfect
Scores in Tests without Studying Much»}: Why Is an Effortless Achiever
Identity Attractive, and for Whom Is It Possible? \emph{Research Papers
in Education}, \emph{30}(4), 393-410.
\url{https://doi.org/10.1080/02671522.2014.970226}

\bibitem[\citeproctext]{ref-jin_meritocracy_2020}
Jin, J., \& Ball, S. J. (2020). Meritocracy, Social Mobility and a New
Form of Class Domination. \emph{British Journal of Sociology of
Education}, \emph{41}(1), 64-79.
\url{https://doi.org/10.1080/01425692.2019.1665496}

\bibitem[\citeproctext]{ref-kelley_subjective_2008}
Kelley, S. M. C., \& Kelley, C. G. E. (2008). \emph{Subjective {Social
Mobility}: {Data} from 30 {Nations}} (\{\{SSRN Scholarly Paper\}\} No.
1306242). Social Science Research Network.
\url{https://papers.ssrn.com/abstract=1306242}

\bibitem[\citeproctext]{ref-khan_saying_2013}
Khan, S., \& Jerolmack, C. (2013). Saying {Meritocracy} and {Doing
Privilege}. \emph{The Sociological Quarterly}, \emph{54}(1), 9-19.
\url{https://doi.org/10.1111/tsq.12008}

\bibitem[\citeproctext]{ref-kluegel_beliefs_1981}
Kluegel, J. R., \& Smith, E. R. (1981). Beliefs about {Stratification}.
\emph{Annu. Rev. Sociol.}, \emph{7}(1), 29-56.
\url{https://doi.org/10.1146/annurev.so.07.080181.000333}

\bibitem[\citeproctext]{ref-kraus_social_2013}
Kraus, M. W., \& Keltner, D. (2013). Social Class Rank, Essentialism,
and Punitive Judgment. \emph{Journal of Personality and Social
Psychology}, \emph{105}(2), 247-261.
\url{https://doi.org/10.1037/a0032895}

\bibitem[\citeproctext]{ref-kwon_access_2026}
Kwon, H., \& Erola, J. (2026). Access to {Higher Education} and
{Support} for {Meritocracy}: {Effort}, {Skill}, and {Education} as
{Components} of {Merit}. \emph{SocScience}, \emph{13}, 997-1017.
\url{https://doi.org/10.15195/v13.a38}

\bibitem[\citeproctext]{ref-lampert_meritocratic_2013}
Lampert, K. (2013). \emph{Meritocratic {Education} and {Social
Worthlessness}}. Palgrave Macmillan UK.
\url{https://doi.org/10.1057/9781137324894}

\bibitem[\citeproctext]{ref-li_community_2021}
Li, A., \& Hu, A. (2021). Community Contexts, Socioeconomic Status, and
Meritocratic Beliefs and Perceptions in {China}. \emph{Chinese Journal
of Sociology}, \emph{7}(3), 470-493.
\url{https://doi.org/10.1177/2057150X211023235}

\bibitem[\citeproctext]{ref-lierse_des_redi_2022}
Lierse, H., Sachweh, P., \& Waitkus, N. (2022). Introduction: {Wealth},
{Inequality} and {Redistribution} in {Capitalist Societies}. \emph{Soc
Just Res}, \emph{35}(4), 367-378.
\url{https://doi.org/10.1007/s11211-022-00402-6}

\bibitem[\citeproctext]{ref-mac_baroz_merit_des_2024}
Mac-Clure, O., Barozet, E., \& Franetovic, G. (2024). {Justicia
distributiva y posición social subjetiva: \textquestiondown la
meritocracia justifica la desigualdad de ingresos?} \emph{Convergencia
Revista de Ciencias Sociales}, \emph{31}, 1-34.
\url{https://doi.org/10.29101/crcs.v31i0.22258}

\bibitem[\citeproctext]{ref-melli_uncovering_2025}
Melli, G., \& Azzollini, L. (2025). Uncovering Misperceptions of Social
Inequalities: What Matters Most, Objective Class or Subjective Social
Status? \emph{Front. Sociol.}, \emph{10}.
\url{https://doi.org/10.3389/fsoc.2025.1617413}

\bibitem[\citeproctext]{ref-merton_social_1938}
Merton, R. K. (1938). Social {Structure} and {Anomie}. \emph{American
Sociological Review}, \emph{3}(5), 672.
\url{https://doi.org/10.2307/2084686}

\bibitem[\citeproctext]{ref-merton_selffulfilling_1948}
Merton, R. K. (1948). The {Self-Fulfilling Prophecy}. \emph{The Antioch
Review}, \emph{8}(2), 193. \url{https://doi.org/10.2307/4609267}

\bibitem[\citeproctext]{ref-meuleman_desert_carin_2020}
Meuleman, B., Roosma, F., \& Abts, K. (2020). Welfare Deservingness
Opinions from Heuristic to Measurable Concept: {The CARIN} Deservingness
Principles Scale. \emph{Social Science Research}, \emph{85}, 102352.
\url{https://doi.org/10.1016/j.ssresearch.2019.102352}

\bibitem[\citeproctext]{ref-mijs_inequality_2018}
Mijs, J. (2018). Inequality {Is} a {Problem} of {Inference}: {How People
Solve} the {Social Puzzle} of {Unequal Outcomes}. \emph{Societies},
\emph{8}(3), 64. \url{https://doi.org/10.3390/soc8030064}

\bibitem[\citeproctext]{ref-mijs_unfulfillable_2016}
Mijs, J. J. B. (2016). The {Unfulfillable Promise} of {Meritocracy}:
{Three Lessons} and {Their Implications} for {Justice} in {Education}.
\emph{Soc Just Res}, \emph{29}(1), 14-34.
\url{https://doi.org/10.1007/s11211-014-0228-0}

\bibitem[\citeproctext]{ref-mijs_paradox_2021}
Mijs, J. J. B. (2021). The Paradox of Inequality: Income Inequality and
Belief in Meritocracy Go Hand in Hand. \emph{Socio-Economic Review},
\emph{19}(1), 7-35. \url{https://doi.org/10.1093/ser/mwy051}

\bibitem[\citeproctext]{ref-mijs_belief_2022}
Mijs, J. J. B., Daenekindt, S., De Koster, W., \& Van Der Waal, J.
(2022). Belief in {Meritocracy Reexamined}: {Scrutinizing} the {Role} of
{Subjective Social Mobility}. \emph{Soc Psychol Q}, \emph{85}(2),
131-141. \url{https://doi.org/10.1177/01902725211063818}

\bibitem[\citeproctext]{ref-mulligan_meritocracy_2023a}
Mulligan, T. (2023). Meritocracy. En E. N. Zalta \& U. Nodelman (Eds.),
\emph{The {Stanford Encyclopedia} of {Philosophy}} (Fall 2023).
Metaphysics Research Lab, Stanford University.
\url{https://plato.stanford.edu/archives/fall2023/entries/meritocracy/}

\bibitem[\citeproctext]{ref-oorschot_blame_2000}
Oorschot, W. V., \& Halman, L. (2000). {BLAME OR FATE}, {INDIVIDUAL OR
SOCIAL}? \emph{European Societies}, \emph{2}(1), 1-28.
\url{https://doi.org/10.1080/146166900360701}

\bibitem[\citeproctext]{ref-parsons_class_1953}
Parsons, T. (1953). A {Revised Analytical Approach} to the {Theory} of
{Social Stratification}. En R. Bendix \& S. M. Lipset (Eds.),
\emph{Class, Status, and Power : A Reader in Social Stratification}.
Glencoe, Ill. : Free Press.
\url{http://archive.org/details/classstatuspower00bendrich}

\bibitem[\citeproctext]{ref-pavel_merit_2024}
Pavel, S. M. (2024). Merit as a {Distributive Justice Principle}. En M.
Sardoč (Ed.), \emph{Handbook of {Equality} of {Opportunity}} (pp.
259-273). Springer International Publishing.
\url{https://doi.org/10.1007/978-3-031-55897-9_93}

\bibitem[\citeproctext]{ref-plaza_networked_2026}
Plaza, A., Beck, G., Iturra-Sanhueza, J., Otero, G., \& Muñoz, B.
(2026). Networked Inequality: {The} Role of Changes in Network
Heterogeneity and Network Size in Attitudes towards Inequality.
\emph{Social Networks}, \emph{84}, 27-45.
\url{https://doi.org/10.1016/j.socnet.2025.07.008}

\bibitem[\citeproctext]{ref-quaresma_merit_u_elites_2024}
Quaresma, M. L., Madrid, S., \& Molina, F. (2024). {Entre el esfuerzo y
el talento: concepciones sobre el mérito de directivos en universidades
de élite y de masas en Chile}. \emph{Revista de la educación superior},
\emph{53}(209), 105-123.
\url{https://doi.org/10.36857/resu.2024.209.2752}

\bibitem[\citeproctext]{ref-quaresma_justify_2024}
Quaresma, M. L., Villalobos, C., Fergnani, M., \& Diaz, B. (2024). I
{Justify Myself}, {Therefore I} Am. {Regimes Of Justification In Chilean
Elite Universities}. \emph{RISE}, 1-19.
\url{https://doi.org/10.17583/rise.7851}

\bibitem[\citeproctext]{ref-renteria_telling_2026}
Rentería, M. (2026). Telling an Exceptional Story: {Meritocratic}
Accounts of Upwardly Mobile {Limeños}. \emph{The Sociological Review},
\emph{74}(6), 1342-1362. \url{https://doi.org/10.1177/00380261251396456}

\bibitem[\citeproctext]{ref-elsoc_data}
Reproducible Research, Centre for Social Conflict and Cohesion Studies
COES. (2024). \emph{Estudio {Longitudinal Social} de {Chile} 2016-2023}
(pp. 5321976, 7154454, 39662924, 179249687, 126465108, 181309930,
179249687, 5332008, 181309930, 7166551, 126607870, 39703925, 7195475,
5332008, 179249687, 181309930, 126679314, 39796123). Harvard Dataverse.
\url{https://doi.org/10.7910/DVN/LD4BPH}

\bibitem[\citeproctext]{ref-retamales_rol_2024}
Retamales, M. (2024). \emph{{El Rol de La Socialización Familiar y
Escolar En La Percepción de Meritocracia En Las Notas: Un Estudio
Multinivel de Estudiantes En Chile}} {[}\{Tesis de Magister\},
Universidad de Chile{]}.
\url{https://jc-castillo.com/references/thesis/posts/2024_el_rol_de_la_sociali}

\bibitem[\citeproctext]{ref-reynolds_perceptions_2014}
Reynolds, J., \& Xian, H. (2014). Perceptions of Meritocracy in the Land
of Opportunity. \emph{Research in Social Stratification and Mobility},
\emph{36}, 121-137. \url{https://doi.org/10.1016/j.rssm.2014.03.001}

\bibitem[\citeproctext]{ref-sandel_tirania_2020}
Sandel, M. J. (2020). \emph{La Tiranía Del Mérito:
\textquestiondown{{Qu\'e}} Ha Sido Del Bien Común?} Debate.

\bibitem[\citeproctext]{ref-sardoc_rise_2024}
Sardoč, M. (2024). The {Rise} of the {Meritocracy}: {The Genealogy} of
an {Idea}. En \emph{Handbook of {Equality} of {Opportunity}} (pp.
329-341). Springer, Cham.
\url{https://doi.org/10.1007/978-3-031-55897-9_7}

\bibitem[\citeproctext]{ref-schrader_school_2015}
Schrader, F.-W., \& Helmke, A. (2015). School {Achievement}:
{Motivational Determinants} and {Processes}. En \emph{International
{Encyclopedia} of the {Social} \& {Behavioral Sciences}} (pp. 48-54).
Elsevier. \url{https://doi.org/10.1016/B978-0-08-097086-8.26055-8}

\bibitem[\citeproctext]{ref-sen_merit_justice_2000}
Sen, A. (2000). Merit and Justice. En \emph{Meritocracy and {Economic
Inequality}}. Princeton University Press.

\bibitem[\citeproctext]{ref-sir_esfuerzo_2017}
Sir, H. (2017). {Más que mérito, esfuerzo. Aproximación crítica al
núcleo ético del mercado del trabajo individualizado en Chile}.
\emph{Trabajo y Sociedad}, \emph{29}, 307-326.
\url{https://www.redalyc.org/articulo.oa?id=387352369016}

\bibitem[\citeproctext]{ref-sonhing_merit_2011}
Son Hing, L. S., Bobocel, D. R., Zanna, M. P., Garcia, D. M., Gee, S.
S., \& Orazietti, K. (2011). The Merit of Meritocracy. \emph{Journal of
Personality and Social Psychology}, \emph{101}(3), 433-450.
\url{https://doi.org/10.1037/a0024618}

\bibitem[\citeproctext]{ref-spruyt_talent_esfuerzo_social_background_2015}
Spruyt, B. (2015). Talent, {Effort} or {Social Background}?
\emph{European Societies}, \emph{17}(1), 94-114.
\url{https://doi.org/10.1080/14616696.2014.977323}

\bibitem[\citeproctext]{ref-tejero-peregrina_perceived_2025}
Tejero-Peregrina, L., Willis, G., Sánchez-Rodríguez, Á., \&
Rodríguez-Bailón, R. (2025). From {Perceived Economic Inequality} to
{Support} for {Redistribution}: {The Role} of {Meritocracy Perception}.
\emph{International Review of Social Psychology}, \emph{38}(1), 4.
\url{https://doi.org/10.5334/irsp.1013}

\bibitem[\citeproctext]{ref-traini_et_al_mov_social_percepciones_meritocracia_2025}
Traini, C., Ehmes, S., \& Gangl, M. (2025). \emph{Climbing the Ladder or
Falling behind: {How} Social Mobility Shapes Perceptions of Meritocracy
in the Wake of Rising Inequality} (Polar \{\{Working\}\} Paper No. 7).
Goethe University.
\url{https://polar-project.org/wp-content/uploads/2025/09/polar_wp7_v250912.pdf}

\bibitem[\citeproctext]{ref-villalobos_percepciones_2022}
Villalobos, C., Kuzmanic, D., Valenzuela, J. P., \& Quaresma, M. L.
(2022). {Percepciones Sociales de los Estudiantes de Educación Superior
de Élite: Un Estudio del Caso Chileno}. \emph{psykhe}, \emph{31}(2).
\url{https://doi.org/10.7764/psykhe.2020.22549}

\bibitem[\citeproctext]{ref-warikoo_legitimating_2014}
Warikoo, N. K., \& Fuhr, C. (2014). Legitimating Status: Perceptions of
Meritocracy and Inequality among Undergraduates at an Elite {British}
University. \emph{British Educational Res J}, \emph{40}(4), 699-717.
\url{https://doi.org/10.1002/berj.3108}

\bibitem[\citeproctext]{ref-weber_etica_2001}
Weber, M. (2001). \emph{{La ética protestante y el {«esp{ı́}ritu»} del
capitalismo}} (J. Abellán, Trad.; 1.ª ed.). Alianza Editorial.

\bibitem[\citeproctext]{ref-weiner_attributional_1985}
Weiner, B. (1985). An Attributional Theory of Achievement Motivation and
Emotion. \emph{Psychological Review}, \emph{92}(4), 548-573.
\url{https://doi.org/10.1037/0033-295X.92.4.548}

\bibitem[\citeproctext]{ref-wiesner_knowing_2026}
Wiesner, T., \& Iturra-Sanhueza, J. (2026). Knowing One's Place? {The}
Role of Income Inequality in Shaping Positioning Bias across 29
Countries. \emph{Front. Sociol.}, \emph{10}.
\url{https://doi.org/10.3389/fsoc.2025.1617014}

\bibitem[\citeproctext]{ref-wiesner_effort_2026}
Wiesner, T., \& Sachweh, P. (2026). Effort versus {Advantage}:
{Visualizing} ({Relative}) {Belief} in {Meritocracy}, 1930 to 2022---{A
Comment} on {Mijs} (2018). \emph{Socius: Sociological Research for a
Dynamic World}, \emph{12}, 23780231261425841.
\url{https://doi.org/10.1177/23780231261425841}

\bibitem[\citeproctext]{ref-wooldridge_aristocracy_2021}
Wooldridge, A. (2021). \emph{The Aristocracy of Talent: How Meritocracy
Made the Modern World}. Skyhorse Publishing.

\bibitem[\citeproctext]{ref-wright_preguntas_2010}
Wright, E. O. (2010). \emph{{Preguntas a la desigualdad: ensayos sobre
análisis de clase, socialismo y marxismo}} (1. ed). Editorial
Universidad del Rosario.

\bibitem[\citeproctext]{ref-young_1964}
Young, M. (1964). \emph{{El triunfo de la meritocracia : 1870-2034 :
ensayo sobre la educación y la igualdad}}. Editorial Tecnos.

\bibitem[\citeproctext]{ref-young_2001}
Young, M. (2001). Down with Meritocracy. \emph{The Guardian}.
\url{https://www.theguardian.com/politics/2001/jun/29/comment}

\bibitem[\citeproctext]{ref-zhu_meritocratic_2025}
Zhu, L. (2025). Meritocratic Beliefs in the {United States}, {Finland},
and {China}: {A} Multidimensional Approach Using Latent Class Analysis.
\emph{British Journal of Sociology}, \emph{76}(1), 153-172.
\url{https://doi.org/10.1111/1468-4446.13152}

\end{CSLReferences}

\cleardoublepage
\phantomsection
\addcontentsline{toc}{part}{Apéndices}
\appendix

\chapter{Instrumentos de
levantamiento}\label{instrumentos-de-levantamiento}

Incluye cuestionarios, guias de entrevista, etc.

\chapter{Tablas adicionales}\label{tablas-adicionales}

Tablas extendidas, tests de robustez, etc.


\backmatter

% --- Cuerpo del libro (capítulos) en arábigos ---
\mainmatter
\pagestyle{scrheadings}

% --- (Opcional) Estilo del índice general ---
% \addtocontents{toc}{\protect\thispagestyle{plain}}


\end{document}
